%% 2-16 何秒後に追いつく？1次方程式（第11回：ページ制限なし。なぜその式・なぜその数を1段ずつ）
\makeatletter
% 式の変形を1行ずつ書く表（左辺｜＝右辺｜この行で使ったきまり）。2-14 の表と同じ形
\NewTblrEnviron{bkiiqhenkei}
\SetTblrInner[bkiiqhenkei]{colspec={Q[r,m] Q[l,m] Q[l,m]}, colsep=0pt, rowsep=1.6pt, stretch=0.85,
  column{2}={rightsep=5mm}, column{3}={leftsep=4mm, font=\sffamily\small, fg=bkDeep},
  vline{3}={0.8pt, bkTheme!45}}
% 上皿てんびん（2-14 の \bkiiobalance と同じ形。つり糸がないので、皿の上の物と線が重ならない）
%   \bkiiqbalance{(x,y)}  皿の上の点は (bkiiqL) (bkiiqR)
\newcommand\bkiiqbalance[1]{\begin{scope}[shift={#1}]
  \fill[bkGray!70, rounded corners=0.6mm] (-13,0) rectangle (13,3.5);
  \fill[bkGray!80] (-1.6,3.5) rectangle (1.6,20);
  \fill[bkInk!75, rounded corners=0.8mm] (-25,20) rectangle (25,22.2);
  \fill[bkRed] (0,21.1) circle (1.5);
  \foreach \s in {-1,1} {
    \fill[bkGray!80] ({\s*21-0.9},22.2) rectangle ({\s*21+0.9},25);
    \fill[bkGray!55, rounded corners=0.5mm] ({\s*21-10},25) rectangle ({\s*21+10},26.8);}
  \coordinate (bkiiqL) at (-21,26.8); \coordinate (bkiiqR) at (21,26.8);
\end{scope}}
% 中の数がわからない箱  \bkiiqbox{(x,y)}  (x,y) は箱の下のまん中。6mm 四方
\newcommand\bkiiqbox[1]{\begin{scope}[shift={#1}]
  \filldraw[fill=bkOrange!25, draw=bkOrange!80!black, line width=0.6pt, rounded corners=0.6mm] (-2.9,0.1) rectangle (2.9,5.9);
  \node[font=\small, inner sep=0pt, text=bkInk] at (0,3) {$t$};
\end{scope}}
% 重さのわかっているおもり  \bkiiqweight{(x,y)}{文字}  (x,y) は下のまん中。はば7mm
\newcommand\bkiiqweight[2]{\begin{scope}[shift={#1}]
  \filldraw[fill=bkGray!40, draw=bkGray!90!black, line width=0.6pt, rounded corners=0.6mm] (-3.5,0.1) rectangle (3.5,5.9);
  \node[font=\sffamily\bfseries\scriptsize, inner sep=0pt, text=bkInk] at (0,3) {#2};
\end{scope}}
% 天びんの下に式を書く  \bkiiqeq{x}{左辺}{右辺}
\newcommand\bkiiqeq[3]{%
  \node[font=\large, inner sep=0pt] at (#1-15.8,-5) {$#2$};
  \node[font=\large, inner sep=0pt] at (#1,-5) {$=$};
  \node[font=\large, inner sep=0pt] at (#1+15.8,-5) {$#3$};}
\makeatother

\bkchapter{2}{16}{\R{何秒後}{なん|びょう|ご}に\R{追}{お}いつく？\\1\R{次方程式}{じ|ほう|てい|しき}}{\bklevel{2}}{\R{何秒後}{なん|びょう|ご}に\R{追}{お}いつく？1\R{次方程式}{じ|ほう|てい|しき}}
\bkmichishirube{2-2,2-12,2-13,2-14}{3-4,3-6,3-10,6-2}

\begin{goal}
\begin{itemize}[label={\color{bkTheme}●}, leftmargin=1.6\zw, itemsep=0pt, topsep=1pt]
\item \bklevel{2}\enspace わからない\R{数}{かず}を\R{文字}{も|じ}にして\R{方程式}{ほう|てい|しき}をつくり、\R{表}{ひょう}でためして\R{解}{かい}をさがせる。
\item \bklevel{2}\enspace \R{等式}{とう|しき}の\R{性質}{せい|しつ}と\R{移項}{い|こう}で、1\R{次方程式}{じ|ほう|てい|しき}を1\R{行}{ぎょう}ずつ\R{解}{と}ける。
\item \bklevel{2}\enspace \R{追}{お}いかけの\R{問題}{もん|だい}を、\R{表}{ひょう}・グラフ・\R{方程式}{ほう|てい|しき}で\R{解}{と}き、\R{代入}{だい|にゅう}して\R{確}{たし}かめられる。
\end{itemize}
\end{goal}

\begin{hatena}
\R{弟}{おとうと}は、\R{兄}{あに}より\qty{10}{m}\R{前}{まえ}から、\R{秒速}{びょう|そく}\qty{2}{m}で\R{走}{はし}り\R{出}{だ}しました。\R{兄}{あに}は\R{同}{おな}じときに、\R{秒速}{びょう|そく}\qty{4}{m}で\R{追}{お}いかけます。\R{兄}{あに}は\R{何秒後}{なん|びょう|ご}に\R{弟}{おとうと}に\R{追}{お}いつくかな？
\begin{bkfigure}
\begin{tikzpicture}[x=1mm, y=1mm]
  \draw[bkGray!60, line width=3pt] (-4,0) -- (118,0);
  \foreach \m in {0,5,...,30} {
    \draw[bkInk, line width=0.7pt] (\m*3.8,-1.6) -- (\m*3.8,1.6);
    \node[bk tick, below=2mm] at (\m*3.8,0) {\m};}
  \node[bk tick, anchor=north west] at (116,-2.5) {m};
  \bkkid{(0,2.5)}{18}{bkBlue}{bkHair}
  \bkkid{(38,2.5)}{14}{bkOrange}{bkHair}
  \draw[bk arrow, bkBlue] (5,12) -- (19,12);
  \node[font=\sffamily\footnotesize, text=bkBlue!80!black, anchor=south] at (12,13.5) {\R{兄}{あに}\enspace\qty{4}{m/s}};
  \draw[bk arrow, bkOrange] (44,10) -- (52,10);
  \node[font=\sffamily\footnotesize, text=bkOrange!80!black, anchor=south] at (50,11.5) {\R{弟}{おとうと}\enspace\qty{2}{m/s}};
  \draw[{Stealth[length=4pt]}-{Stealth[length=4pt]}, bkRed, line width=0.7pt] (1,-9.5) -- (37,-9.5);
  \node[font=\sffamily\footnotesize, text=bkRed!80!black, anchor=north] at (19,-10.5) {\qty{10}{m}\R{先}{さき}};
  \node[font=\sffamily\bfseries\Large, text=bkOrange!80!black] at (90,12) {?};
\end{tikzpicture}
\bkfigcaption{2\R{人}{にん}は\R{同}{おな}じときにスタート。どこで\R{追}{お}いつく？}
\end{bkfigure}
\end{hatena}

\noindent \R{兄}{あに}のほうが\R{速}{はや}いので、いつかは\R{追}{お}いつきそうです。でも、それが\R{何秒後}{なん|びょう|ご}なのかは、\R{見}{み}ているだけではわかりません。この\R{章}{しょう}では、「わからない\R{数}{かず}」を\R{文字}{も|じ}にして\R{式}{しき}をつくり、その\R{数}{かず}を\R{見}{み}つけ\R{出}{だ}す\R{方法}{ほう|ほう}を\R{学}{まな}びます。\R{同}{おな}じ\R{問題}{もん|だい}を、\R{表}{ひょう}・グラフ・\R{方程式}{ほう|てい|しき}の3つの\R{方法}{ほう|ほう}で\R{解}{と}いてくらべます。

%% ---------------- やってみよう・思い出そう ----------------
\begin{yattemiyou}[\R{電卓}{でん|たく}でためして、\R{表}{ひょう}を\R{作}{つく}ろう]
スタートから $t$ \R{秒後}{びょう|ご}の2\R{人}{にん}の\R{位置}{い|ち}を\R{計算}{けい|さん}します。\R{位置}{い|ち}は、スタートのときに\R{兄}{あに}がいた\R{場所}{ば|しょ}を0として、そこから\R{前}{まえ}へ\R{何}{なん}mの\R{所}{ところ}にいるかで\R{表}{あらわ}します。
\begin{enumerate}[label={\bkcircled[bkGreen]{\arabic*}}, leftmargin=2.4\zw, itemsep=1pt, topsep=2pt]
\item \R{兄}{あに}は1\R{秒}{びょう}に\qty{4}{m}、\R{弟}{おとうと}は1\R{秒}{びょう}に\qty{2}{m}\R{進}{すす}みます。\R{下}{した}の\R{表}{ひょう}の1\R{秒後}{びょう|ご}・2\R{秒後}{びょう|ご}の\R{数}{かず}が、なぜそうなるか\R{確}{たし}かめよう。
\item \R{表}{ひょう}をノートに\R{写}{うつ}し、$t=3$ から\R{先}{さき}を\R{電卓}{でん|たく}で\R{計算}{けい|さん}してうめよう。
\item 2\R{人}{にん}の\R{位置}{い|ち}が\R{同}{おな}じになるのは、\R{何秒後}{なん|びょう|ご}かな？
\end{enumerate}
\begin{bkcenter}
\begin{bktable}{colspec={Q[l,m,wd=33mm] *{8}{Q[c,m,wd=7.2mm]}}, rowsep=1pt, colsep=3pt}
\R{時間}{じ|かん} $t$（s） & 0 & 1 & 2 & 3 & 4 & 5 & 6 & 7 \\
\R{兄}{あに}の\R{位置}{い|ち}（m） & 0 & 4 & 8 & & & & & \\
\R{弟}{おとうと}の\R{位置}{い|ち}（m） & 10 & 12 & 14 & & & & & \\
\end{bktable}
\end{bkcenter}
\end{yattemiyou}

\begin{omoidasou}
\bkrecall{2-3}{\R{道}{みち}のり＝\R{速}{はや}さ×\R{時間}{じ|かん}。\qty{4}{m/s}で3\R{秒}{びょう}\R{進}{すす}むと $4\times3=12$ で\qty{12}{m}。}
\bkrecall{2-12}{$4\times t$ は×をはぶいて $4t$ と\R{書}{か}く。\R{同類項}{どう|るい|こう}をまとめる：$4t-2t=2t$。\R{分配法則}{ぶん|ぱい|ほう|そく}：$3(x+4)=3x+12$}
\bkrecall{2-14}{\R{等式}{とう|しき}は\R{天}{てん}びん。\R{両辺}{りょう|へん}に\R{同}{おな}じことをしても、\R{等式}{とう|しき}は\R{成}{な}り\R{立}{た}つ。\R{移項}{い|こう}すると\R{符号}{ふ|ごう}が\R{反対}{はん|たい}になる。}
\bkrecall{2-2}{$(-10)\div(-2)=5$。マイナスどうしのわり\R{算}{ざん}の\R{答}{こた}えはプラス。}
\bkrecall{2-10}{1\R{分}{ぷん}あたり60\R{秒}{びょう}なので \qty{60}{s/min}。$\qty{2}{min}\times\qty{60}{s/min}=\qty{120}{s}$（minが\R{約分}{やく|ぶん}されて\R{消}{き}える）}
\end{omoidasou}



\clearpage
%% ---------------- 数学の道具 ----------------
\bksection{doguu}{\R{数学}{すう|がく}の\R{道具}{どう|ぐ}}

\bksubsection{1}{\R{位置}{い|ち}を\R{式}{しき}で\R{表}{あらわ}す：なぜ $4t$ と $10+2t$}

\leavevmode\label{bk:2-16:ichi}まず、2\R{人}{にん}が\R{何秒後}{なん|びょう|ご}にどこにいるかを、\R{図}{ず}にしてみましょう。\R{青}{あお}い\R{点}{てん}が\R{兄}{あに}、だいだいの\R{点}{てん}が\R{弟}{おとうと}です。1\R{秒}{びょう}ごとに、\R{兄}{あに}は\qty{4}{m}、\R{弟}{おとうと}は\qty{2}{m}だけ\R{右}{みぎ}へ\R{進}{すす}みます。

\begin{bkfigure}
\begin{tikzpicture}[x=1mm, y=1mm]
  \foreach \m in {0,4,...,24} \draw[bkLine, line width=0.4pt, dashed] (26+\m*4,-1.5) -- (26+\m*4,75);
  \foreach \t in {0,...,5} {
    \pgfmathsetmacro\yy{70-\t*14}
    \pgfmathsetmacro\xa{26+4*\t*4}
    \pgfmathsetmacro\xb{26+(10+2*\t)*4}
    \pgfmathtruncatemacro\gap{10-2*\t}
    \draw[bkGray!70, line width=1.2pt] (26,\yy) -- (124,\yy);
    \node[anchor=west, font=\sffamily\small, text=bkInk] at (0,\yy) {\t\,\R{秒後}{びょう|ご}};
    \ifnum\gap>0
      \draw[{Stealth[length=3pt]}-{Stealth[length=3pt]}, bkRed, line width=0.6pt] (\xa+1.8,\yy+3.6) -- (\xb-1.8,\yy+3.6);
      \node[anchor=south, font=\sffamily\footnotesize, text=bkRed!80!black, inner sep=0.8pt] at ({(\xa+\xb)/2},\yy+4.4) {\R{差}{さ} \gap\,m};
      \fill[bkBlue] (\xa,\yy) circle (1.7);
      \fill[bkOrange] (\xb,\yy) circle (1.7);
    \else
      \fill[bkBlue] (\xa-1.2,\yy) circle (1.7);
      \fill[bkOrange] (\xb+1.2,\yy) circle (1.7);
      \node[anchor=south, font=\sffamily\bfseries\footnotesize, text=bkRed!80!black, inner sep=0.8pt] at (\xa,\yy+3) {\R{同}{おな}じ\R{場所}{ば|しょ}！};
    \fi}
  \foreach \m in {0,4,...,24} \node[bk tick, anchor=north] at (26+\m*4,-2.5) {\m};
  \node[bk tick, anchor=north west] at (124,-2.5) {m};
  \fill[bkBlue] (40,-12) circle (1.7); \node[anchor=west, font=\sffamily\footnotesize] at (42.5,-12) {\R{兄}{あに}（1\R{秒}{びょう}に\qty{4}{m}）};
  \fill[bkOrange] (82,-12) circle (1.7); \node[anchor=west, font=\sffamily\footnotesize] at (84.5,-12) {\R{弟}{おとうと}（1\R{秒}{びょう}に\qty{2}{m}）};
\end{tikzpicture}
\bkfigcaption{1\R{秒}{びょう}ごとの2\R{人}{にん}の\R{位置}{い|ち}。\R{差}{さ}（\R{赤}{あか}）は、1\R{秒}{びょう}に\qty{2}{m}ずつ\R{小}{ちい}さくなる。}
\end{bkfigure}

\textbf{\R{兄}{あに}の\R{位置}{い|ち}}\quad \R{兄}{あに}は0の\R{所}{ところ}から\R{出}{で}て、1\R{秒}{びょう}に\qty{4}{m}ずつ\R{進}{すす}みます。1\R{秒後}{びょう|ご}は $4\times1=4$、2\R{秒後}{びょう|ご}は $4\times2=8$、3\R{秒後}{びょう|ご}は $4\times3=12$ です。どれも「4×\R{秒}{びょう}の\R{数}{かず}」になっています。これは「\R{道}{みち}のり＝\R{速}{はや}さ×\R{時間}{じ|かん}」（2-3）そのものです。\R{秒}{びょう}の\R{数}{かず}を $t$ とすると、\R{兄}{あに}の\R{位置}{い|ち}は $4\times t$、×をはぶいて $4t$（m）です（2-12）。

\textbf{\R{弟}{おとうと}の\R{位置}{い|ち}}\quad \R{弟}{おとうと}は、はじめから10の\R{所}{ところ}にいます。そこから1\R{秒}{びょう}に\qty{2}{m}ずつ\R{進}{すす}むので、$t$ \R{秒}{びょう}で $2t$（m）\R{進}{すす}みます。\R{位置}{い|ち}は「はじめの10」に「\R{進}{すす}んだ $2t$」をたした $10+2t$（m）です。3\R{秒後}{びょう|ご}なら $10+2\times3=16$ です。×を＋より\R{先}{さき}に\R{計算}{けい|さん}します（2-11）。

\begin{bkcenter}
\begin{bktable}{colspec={Q[l,m,wd=30mm] *{6}{Q[c,m,wd=12.5mm]}}, rowsep=1pt, colsep=2.5pt, cells={font=\small}, column{1}={font=\sffamily\small}}
\R{時間}{じ|かん} $t$（s） & 0 & 1 & 2 & 3 & 4 & 5 \\
\R{兄}{あに} $4t$（m） & {$4\times0$\\$=0$} & {$4\times1$\\$=4$} & {$4\times2$\\$=8$} & {$4\times3$\\$=12$} & {$4\times4$\\$=16$} & {$4\times5$\\$=20$} \\
\R{弟}{おとうと} $10+2t$（m） & {$10+0$\\$=10$} & {$10+2$\\$=12$} & {$10+4$\\$=14$} & {$10+6$\\$=16$} & {$10+8$\\$=18$} & {$10+10$\\$=20$} \\
\end{bktable}
\bktabcaption{\R{式}{しき}の $t$ に0、1、2、…を\R{代入}{だい|にゅう}すると（2-13）、やってみようの\R{表}{ひょう}と\R{同}{おな}じ\R{数}{かず}が\R{出}{で}る。}
\end{bkcenter}

\bksubsection{2}{\R{差}{さ}に\R{目}{め}をつける：1\R{秒}{びょう}に\qty{2}{m}ずつ\R{近}{ちか}づく}

\leavevmode\label{bk:2-16:sa}もう1つの\R{見方}{み|かた}もあります。\R{上}{うえ}の\R{図}{ず}の\R{赤}{あか}い\R{差}{さ}を\R{見}{み}てください。10、8、6、4、2、0と、1\R{秒}{びょう}に2ずつへっています。

なぜ2なのでしょう。1\R{秒}{びょう}の\R{間}{あいだ}に、\R{兄}{あに}は\qty{4}{m}\R{進}{すす}み、\R{弟}{おとうと}は\qty{2}{m}\R{進}{すす}みます。\R{兄}{あに}のほうが $4-2=2$ で\qty{2}{m}\R{多}{おお}く\R{進}{すす}むので、そのぶん\R{差}{さ}がちぢまるのです。\R{兄}{あに}は\R{弟}{おとうと}に、1\R{秒}{びょう}に\qty{2}{m}ずつ\R{近}{ちか}づいていきます。これを「\R{近}{ちか}づく\R{速}{はや}さ」とよぶことにします。
\[ \qty{4}{m/s}-\qty{2}{m/s}=\qty{2}{m/s} \]
はじめの\R{差}{さ}は\qty{10}{m}です。\qty{2}{m}ずつ\R{近}{ちか}づいて\qty{10}{m}をうめるには、10の\R{中}{なか}に2が\R{何}{なん}こあるかを\R{考}{かんが}えて、$10\div2=5$ で5\R{秒}{びょう}かかります。\R{単位}{たん|い}ごと\R{書}{か}くと（2-13）、
\[ \qty{10}{m}\div\qty{2}{m/s}=10\,\mathrm{m}\times\frac{1\,\mathrm{s}}{2\,\mathrm{m}}=\qty{5}{s} \]
mどうしが\R{約分}{やく|ぶん}されて、\R{時間}{じ|かん}の\R{単位}{たん|い}sが\R{残}{のこ}ります。\R{答}{こた}えは5\R{秒}{びょう}。\R{表}{ひょう}でも、5\R{秒後}{びょう|ご}に2\R{人}{にん}とも\qty{20}{m}の\R{所}{ところ}にいました。

\begin{bkcenter}
\begin{bktable}{colspec={Q[l,m,wd=36mm] *{6}{Q[c,m,wd=9mm]}}, rowsep=1pt, colsep=3pt, column{7}={bg=bkYellow!20}}
\R{時間}{じ|かん} $t$（s） & 0 & 1 & 2 & 3 & 4 & 5 \\
\R{差}{さ}（\R{弟}{おとうと}－\R{兄}{あに}）（m） & 10 & 8 & 6 & 4 & 2 & 0 \\
\end{bktable}
\bktabcaption{\R{差}{さ}は1\R{秒}{びょう}に2ずつへり、5\R{秒後}{びょう|ご}に0になる。}
\end{bkcenter}

\bkhatenaref の\R{答}{こた}え：\R{兄}{あに}は\textbf{5\R{秒後}{びょう|ご}}に、スタートから\qty{20}{m}の\R{所}{ところ}で\R{弟}{おとうと}に\R{追}{お}いつきます。

\R{差}{さ}を\R{使}{つか}うこの\R{考}{かんが}え\R{方}{かた}は、\R{追}{お}いかけの\R{問題}{もん|だい}にはとても\R{便利}{べん|り}です。でも、\R{水}{みず}そうの\R{水}{みず}がふえる\R{問題}{もん|だい}や、ものを\R{配}{くば}る\R{問題}{もん|だい}など、\R{場面}{ば|めん}が\R{変}{か}わると、そのたびに\R{新}{あたら}しい\R{考}{かんが}え\R{方}{かた}を\R{工夫}{く|ふう}しなければなりません。そこで、どんな\R{場面}{ば|めん}でも\R{同}{おな}じ\R{手順}{て|じゅん}で\R{解}{と}ける\R{方法}{ほう|ほう}を、\R{次}{つぎ}から\R{学}{まな}びます。

\bksubsection{3}{ためして\R{探}{さが}す：\R{方程式}{ほう|てい|しき}と\R{解}{かい}}

\leavevmode\label{bk:2-16:houteishiki}「\R{追}{お}いつく」とは、「2\R{人}{にん}の\R{位置}{い|ち}が\R{等}{ひと}しくなる」ことです。\R{兄}{あに}の\R{位置}{い|ち}は $4t$、\R{弟}{おとうと}の\R{位置}{い|ち}は $10+2t$ だったので、「＝」でつなぐと\R{等式}{とう|しき}（2-14）になります。
\[ 4t=10+2t \]
\R{位置}{い|ち}の\R{単位}{たん|い}はm、\R{時間}{じ|かん}の\R{単位}{たん|い}はsにそろえてあるので、\R{式}{しき}の\R{中}{なか}では\R{数}{かず}だけを\R{書}{か}いています。

この\R{等式}{とう|しき}の $t$ に、いろいろな\R{数}{かず}を\R{入}{い}れて、\R{左辺}{さ|へん}と\R{右辺}{う|へん}をくらべてみましょう。

\begin{bkcenter}
\begin{bktable}{colspec={Q[c,m,wd=12mm] Q[c,m,wd=30mm] Q[c,m,wd=36mm] Q[c,m,wd=26mm]}, rowsep=1pt, colsep=3pt, cells={font=\small}, row{1}={font=\sffamily\bfseries\small}, row{7}={bg=bkYellow!20}}
$t$ & \R{左辺}{さ|へん} $4t$ & \R{右辺}{う|へん} $10+2t$ & \R{等}{ひと}しい？ \\
1 & $4\times1=4$ & $10+2\times1=12$ & ×（$4<12$） \\
2 & $4\times2=8$ & $10+2\times2=14$ & ×（$8<14$） \\
3 & $4\times3=12$ & $10+2\times3=16$ & ×（$12<16$） \\
4 & $4\times4=16$ & $10+2\times4=18$ & ×（$16<18$） \\
5 & $4\times5=20$ & $10+2\times5=20$ & ○（$20=20$） \\
6 & $4\times6=24$ & $10+2\times6=22$ & ×（$24>22$） \\
\end{bktable}
\bktabcaption{$t=5$ のときだけ、\R{左辺}{さ|へん}と\R{右辺}{う|へん}が\R{等}{ひと}しくなる。}
\end{bkcenter}

$t=3$ のときは $12=16$ となり、まちがった\R{等式}{とう|しき}です。$t=5$ のときだけ、$20=20$ で\R{正}{ただ}しい\R{等式}{とう|しき}になります。6\R{秒後}{びょう|ご}には\R{兄}{あに}のほうが\R{前}{まえ}に\R{出}{で}るので、また\R{等}{ひと}しくなくなります。

このように、\R{文字}{も|じ}に\R{入}{い}れる\R{数}{かず}によって、\R{成}{な}り\R{立}{た}ったり\R{成}{な}り\R{立}{た}たなかったりする\R{等式}{とう|しき}を\textbf{\R{方程式}{ほう|てい|しき}}といいます。\bkidx{ほうていしき}{\R{方程式}{ほう|てい|しき}}\R{方程式}{ほう|てい|しき}を\R{成}{な}り\R{立}{た}たせる\R{数}{かず}（ここでは5）を、その\R{方程式}{ほう|てい|しき}の\textbf{\R{解}{かい}}といいます。\R{解}{かい}を\R{求}{もと}めることを\textbf{\R{方程式}{ほう|てい|しき}を\R{解}{と}く}といいます。\bkidx{かい}{\R{解}{かい}（\R{方程式}{ほう|てい|しき}の）}

\begin{bkcenter}
\begin{bktable}{colspec={Q[l,m,wd=40mm] Q[l,m,wd=78mm]}, rowsep=1pt, colsep=4pt, cells={font=\small}, row{1}={font=\sffamily\bfseries\small}}
\R{式}{しき}の\R{種類}{しゅ|るい} & \R{例}{れい}と、\R{文字}{も|じ}に\R{数}{かず}を\R{入}{い}れたとき \\
いつでも\R{成}{な}り\R{立}{た}つ\R{等式}{とう|しき} & $a+b=b+a$（2-11）。どんな\R{数}{かず}を\R{入}{い}れても\R{成}{な}り\R{立}{た}つ \\
\R{方程式}{ほう|てい|しき} & $4t=10+2t$。$t=5$ のときだけ\R{成}{な}り\R{立}{た}つ \\
\end{bktable}
\bktabcaption{「＝」のある\R{式}{しき}でも、いつでも\R{成}{な}り\R{立}{た}つものと、\R{方程式}{ほう|てい|しき}がある。}
\end{bkcenter}

\R{表}{ひょう}で1つずつためすのは、\R{確}{たし}かな\R{方法}{ほう|ほう}です。でも、\R{答}{こた}えが100\R{秒後}{びょう|ご}なら100\R{回}{かい}も\R{計算}{けい|さん}しなければなりません。もっと\R{早}{はや}く\R{解}{と}く\R{方法}{ほう|ほう}を\R{考}{かんが}えましょう。

\bksubsection{4}{\R{天}{てん}びんで\R{解}{と}く}

\leavevmode\label{bk:2-16:tenbin}2-14で、\R{等式}{とう|しき}は\R{天}{てん}びんのつり\R{合}{あ}いと\R{同}{おな}じだと\R{学}{まな}びました。\R{左右}{さ|ゆう}の\R{皿}{さら}から\R{同}{おな}じ\R{重}{おも}さをとっても、\R{同}{おな}じ\R{重}{おも}さを\R{分}{わ}けても、つり\R{合}{あ}ったままです。

$4t=10+2t$ を\R{天}{てん}びんにしてみましょう。\R{中}{なか}の\R{数}{かず}がわからない\R{箱}{はこ}を $t$ とします。$4t$ は「$t$ の\R{箱}{はこ}が4つ」、$10+2t$ は「10のおもりと $t$ の\R{箱}{はこ}2つ」です。

\begin{bkfigure}
\begin{tikzpicture}[x=1mm, y=1mm]
\begin{scope}[scale=0.75]
  \bkiiqbalance{(0,0)}
  \bkiiqbox{($(bkiiqL)+(-3,0)$)}\bkiiqbox{($(bkiiqL)+(3,0)$)}\bkiiqbox{($(bkiiqL)+(-3,6)$)}\bkiiqbox{($(bkiiqL)+(3,6)$)}
  \bkiiqweight{($(bkiiqR)+(-6.5,0)$)}{10}\bkiiqbox{($(bkiiqR)+(0.5,0)$)}\bkiiqbox{($(bkiiqR)+(6.5,0)$)}
\end{scope}
\begin{scope}[shift={(82,0)}, scale=0.75]
  \bkiiqbalance{(0,0)}
  \bkiiqbox{($(bkiiqL)+(-3,0)$)}\bkiiqbox{($(bkiiqL)+(3,0)$)}
  \bkiiqweight{(bkiiqR)}{10}
\end{scope}
  \bkiiqeq{0}{4t}{10+2t}
  \bkiiqeq{82}{2t}{10}
  \draw[bk arrow, bkTheme] (30,10) -- (52,10);
  \node[font=\sffamily\footnotesize, text=bkInk, align=center, anchor=south] at (41,12) {\R{両方}{りょう|ほう}から\\\R{箱}{はこ}を2つずつとる};
\end{tikzpicture}
\bkfigcaption{①\enspace \R{両方}{りょう|ほう}の\R{皿}{さら}から\R{箱}{はこ}を2つずつとると、\R{箱}{はこ}2つが10とつり\R{合}{あ}う。}
\end{bkfigure}

\textbf{なぜ2つずつとるの？}\quad \R{右}{みぎ}の\R{皿}{さら}には、\R{箱}{はこ}が2つとおもり10がのっています。\R{箱}{はこ}を2つとれば、\R{右}{みぎ}の\R{皿}{さら}はおもり10だけになり、\R{箱}{はこ}は\R{左}{ひだり}の\R{皿}{さら}だけに\R{残}{のこ}ります。\R{箱}{はこ}を\R{片方}{かた|ほう}に\R{集}{あつ}めると、\R{箱}{はこ}の\R{重}{おも}さを\R{考}{かんが}えやすくなります。\R{左}{ひだり}の\R{皿}{さら}からも\R{同}{おな}じ2つをとるので、つり\R{合}{あ}いはくずれません。

\begin{bkfigure}
\begin{tikzpicture}[x=1mm, y=1mm]
\begin{scope}[scale=0.75]
  \bkiiqbalance{(0,0)}
  \bkiiqbox{($(bkiiqL)+(-3,0)$)}\bkiiqbox{($(bkiiqL)+(3,0)$)}
  \bkiiqweight{(bkiiqR)}{10}
\end{scope}
\begin{scope}[shift={(82,0)}, scale=0.75]
  \bkiiqbalance{(0,0)}
  \bkiiqbox{(bkiiqL)}
  \bkiiqweight{(bkiiqR)}{5}
\end{scope}
  \bkiiqeq{0}{2t}{10}
  \bkiiqeq{82}{t}{5}
  \draw[bk arrow, bkTheme] (30,10) -- (52,10);
  \node[font=\sffamily\footnotesize, text=bkInk, align=center, anchor=south] at (41,12) {\R{両方}{りょう|ほう}を\\2つに\R{分}{わ}ける};
\end{tikzpicture}
\bkfigcaption{②\enspace \R{箱}{はこ}2つで10なら、\R{両方}{りょう|ほう}を\R{半分}{はん|ぶん}にすると、\R{箱}{はこ}1つは5とつり\R{合}{あ}う。}
\end{bkfigure}

\textbf{なぜ2でわるの？}\quad \R{箱}{はこ}2つで10の\R{重}{おも}さです。\R{同}{おな}じ\R{箱}{はこ}が2つなので、1つ\R{分}{ぶん}は10を2つに\R{等}{ひと}しく\R{分}{わ}けた $10\div2=5$ です。\R{両方}{りょう|ほう}の\R{皿}{さら}を\R{半分}{はん|ぶん}にしても、つり\R{合}{あ}いはくずれません。\R{式}{しき}では「\R{両辺}{りょう|へん}を2でわる」です（2-14）。

\R{天}{てん}びんでしたことを、\R{式}{しき}で1\R{行}{ぎょう}ずつ\R{書}{か}くと、こうなります。\R{右}{みぎ}のらんには、その\R{行}{ぎょう}でしたことを\R{書}{か}きます。

\begin{bkcenter}
\begin{bkiiqhenkei}{}
$4t$ & ${}=10+2t$ & \\
$4t-2t$ & ${}=10+2t-2t$ & ①\R{両辺}{りょう|へん}から $2t$ をひく（\R{箱}{はこ}を2つずつとる） \\
$2t$ & ${}=10$ & \R{同類項}{どう|るい|こう}をまとめる（$4t-2t=2t$、$2t-2t=0$） \\
$t$ & ${}=5$ & ②\R{両辺}{りょう|へん}を2でわる（2つに\R{分}{わ}ける） \\
\end{bkiiqhenkei}
\end{bkcenter}

\R{答}{こた}えは $t=5$ で、\R{表}{ひょう}でためしたときと\R{同}{おな}じです。\R{表}{ひょう}なら6\R{回}{かい}も\R{計算}{けい|さん}しましたが、\R{天}{てん}びんの\R{考}{かんが}え\R{方}{かた}なら2つの\R{操作}{そう|さ}で\R{答}{こた}えにたどりつきました。

\bksubsection{5}{\R{移項}{い|こう}で\R{短}{みじか}く\R{書}{か}く}

\leavevmode\label{bk:2-16:ikou}\R{上}{うえ}の2\R{行目}{ぎょう|め}では、\R{両辺}{りょう|へん}から $2t$ をひきました。\R{右辺}{う|へん}では $+2t-2t$ で $2t$ が\R{消}{き}え、\R{左辺}{さ|へん}には $-2t$ がふえました。\R{結果}{けっ|か}だけを\R{見}{み}ると、\R{右辺}{う|へん}の $+2t$ が、\R{符号}{ふ|ごう}を\R{変}{か}えて $-2t$ として\R{左辺}{さ|へん}へ\R{移}{うつ}ったように\R{見}{み}えます。

\begin{bkfigure}
\begin{tikzpicture}[x=1mm, y=1mm]
  \node[font=\large, inner sep=1pt] (l1) at (0,12) {$4t$};
  \node[font=\large, inner sep=1pt] at (10,12) {$=$};
  \node[font=\large, inner sep=1pt] at (20,12) {$10$};
  \node[font=\large, inner sep=1pt, text=bkBlue!80!black] (r1) at (32,12) {$+2t$};
  \node[font=\large, inner sep=1pt] at (-6,0) {$4t$};
  \node[font=\large, inner sep=1pt, text=bkBlue!80!black] (l2) at (6,0) {$-2t$};
  \node[font=\large, inner sep=1pt] at (17,0) {$=$};
  \node[font=\large, inner sep=1pt] at (26,0) {$10$};
  \draw[bkBlue, line width=0.8pt, -{Stealth[length=4pt]}, shorten <=0.5mm, shorten >=0.5mm] (r1.south) -- (l2.north east);
  \node[anchor=west, font=\sffamily\footnotesize, text=bkInk, align=left] at (44,6) {「＝」の\R{向}{む}こうへうつすと、\\\R{符号}{ふ|ごう}が\R{反対}{はん|たい}（＋→－）になる。\\\R{中身}{なか|み}は「\R{両辺}{りょう|へん}から $2t$ をひく」こと。};
\end{tikzpicture}
\bkfigcaption{\R{移項}{い|こう}（2-14）。\R{両辺}{りょう|へん}から\R{同}{おな}じものをひく\R{計算}{けい|さん}を、1\R{行}{ぎょう}で\R{書}{か}いたもの。}
\end{bkfigure}

このように\R{項}{こう}を\R{反対}{はん|たい}の\R{辺}{へん}へうつすことが、2-14で\R{学}{まな}んだ\textbf{\R{移項}{い|こう}}です。\R{移項}{い|こう}を\R{使}{つか}うと、「\R{両辺}{りょう|へん}から $2t$ をひく」の\R{行}{ぎょう}を\R{書}{か}かずにすみます。
\[ 4t=10+2t\quad\to\quad 4t-2t=10 \]
マイナスの\R{項}{こう}を\R{移項}{い|こう}すると、プラスになります。たとえば $x-7=5$ の $-7$ を\R{移項}{い|こう}すると $x=5+7$ です。これは「\R{両辺}{りょう|へん}に7をたす」ことを1\R{行}{ぎょう}で\R{書}{か}いたものです（$x-7+7=5+7$）。

\bksubsection{6}{1\R{次方程式}{じ|ほう|てい|しき}の\R{解}{と}き\R{方}{かた}}

\leavevmode\label{bk:2-16:tejun}$4t=10+2t$ の\R{文字}{も|じ}の\R{項}{こう}は $4t$ と $2t$ です。どちらも「\R{数}{かず}×\R{文字}{も|じ}」で、\R{文字}{も|じ}を1\R{回}{かい}だけかけています。$t^2$（$t\times t$、2-7）のように\R{文字}{も|じ}を2\R{回}{かい}かけた\R{項}{こう}はありません。このように、\R{文字}{も|じ}の\R{項}{こう}が「\R{数}{かず}×\R{文字}{も|じ}」だけの\R{方程式}{ほう|てい|しき}を\textbf{1\R{次方程式}{じ|ほう|てい|しき}}といいます。\bkidx{いちじほうていしき}{1\R{次方程式}{じ|ほう|てい|しき}}「1\R{次}{じ}」は、\R{文字}{も|じ}を1\R{回}{かい}かけている、という\R{意味}{い|み}です（$t^2$ のある\R{方程式}{ほう|てい|しき}は2\R{次方程式}{じ|ほう|てい|しき}で、3-10で\R{学}{まな}びます）。

1\R{次方程式}{じ|ほう|てい|しき}は、いつも\R{次}{つぎ}の3つの\R{手順}{て|じゅん}で\R{解}{と}けます。

\begin{bkcenter}
\begin{bktable}{colspec={Q[l,m,wd=58mm] Q[l,m,wd=62mm]}, rowsep=2pt, colsep=4pt, cells={font=\small}, row{1}={font=\sffamily\bfseries\small}}
\R{手順}{て|じゅん} & なぜそうするの？ \\
① \R{文字}{も|じ}の\R{項}{こう}を\R{左辺}{さ|へん}に、\R{数}{かず}の\R{項}{こう}を\R{右辺}{う|へん}に\R{移項}{い|こう}する & \R{天}{てん}びんで、\R{箱}{はこ}を\R{片方}{かた|ほう}の\R{皿}{さら}に、おもりをもう\R{片方}{かた|ほう}に\R{集}{あつ}めるため \\
② \R{両辺}{りょう|へん}をそれぞれまとめて、$ax=b$ の\R{形}{かたち}にする & 「\R{箱}{はこ}が $a$ こで、おもり $b$ とつり\R{合}{あ}う」\R{形}{かたち}にするため \\
③ \R{両辺}{りょう|へん}を $x$ の\R{係数}{けい|すう} $a$ でわる & \R{箱}{はこ} $a$ こ\R{分}{ぶん}から、\R{箱}{はこ}1こ\R{分}{ぶん}を\R{求}{もと}めるため \\
\end{bktable}
\bktabcaption{1\R{次方程式}{じ|ほう|てい|しき}の\R{解}{と}き\R{方}{かた}。\R{天}{てん}びんでしたことと\R{同}{おな}じ。}
\end{bkcenter}

たとえば $3x-7=x+5$ を\R{解}{と}いてみましょう。\R{文字}{も|じ}は $t$ でも $x$ でも、\R{解}{と}き\R{方}{かた}は\R{同}{おな}じです。

\begin{bkcenter}
\begin{bkiiqhenkei}{}
$3x-7$ & ${}=x+5$ & \\
$3x-x$ & ${}=5+7$ & ①\R{右辺}{う|へん}の $x$ を\R{左辺}{さ|へん}へ、\R{左辺}{さ|へん}の $-7$ を\R{右辺}{う|へん}へ\R{移項}{い|こう} \\
$2x$ & ${}=12$ & ②まとめる（$3x-x=3x-1x=2x$） \\
$x$ & ${}=6$ & ③\R{両辺}{りょう|へん}を2でわる（$12\div2=6$） \\
\end{bkiiqhenkei}
\end{bkcenter}

\begin{itemize}[leftmargin=1.2\zw, itemsep=1pt, topsep=2pt, label={\color{bkTheme}●}]
\item $x$ は $1x$ のことです（2-12）。だから $3x-x$ は $3x-1x=2x$ です。
\item $-7$ は\R{移項}{い|こう}すると $+7$ になります。\R{両辺}{りょう|へん}に7をたしたのと\R{同}{おな}じです。
\item $+x$ は\R{移項}{い|こう}すると $-x$ になります。\R{両辺}{りょう|へん}から $x$ をひいたのと\R{同}{おな}じです。
\end{itemize}

\bksubsection{7}{\R{表}{ひょう}・グラフ・\R{方程式}{ほう|てい|しき}で、\R{同}{おな}じ\R{答}{こた}え}

\leavevmode\label{bk:2-16:graph}3つ\R{目}{め}の\R{方法}{ほう|ほう}は、グラフです。\R{横}{よこ}に\R{時間}{じ|かん} $t$、たてに\R{位置}{い|ち}をとります（\R{比例}{ひ|れい}のグラフ、2-15）。グラフは、\R{次}{つぎ}の\R{手順}{て|じゅん}でかきます。

\begin{enumerate}[label={\bkcircled{\arabic*}}, leftmargin=2.4\zw, itemsep=0pt, topsep=2pt]
\item \R{表}{ひょう}を\R{作}{つく}る（\R{兄}{あに}は $4t$、\R{弟}{おとうと}は $10+2t$ に $t=0$、1、2、…を\R{代入}{だい|にゅう}）。
\item \R{表}{ひょう}の\R{数}{かず}の\R{組}{く}を、\R{点}{てん}にしてとる（$t=1$、\R{位置}{い|ち}4 なら、\R{横}{よこ}1・たて4の\R{所}{ところ}）。
\item \R{同}{おな}じ\R{人}{ひと}の\R{点}{てん}を、まっすぐな\R{線}{せん}でつなぐ（\R{速}{はや}さが\R{一定}{いっ|てい}なので、\R{点}{てん}は\R{一直線}{いっ|ちょく|せん}にならぶ）。
\item 2\R{本}{ほん}の\R{線}{せん}が\R{交}{まじ}わる\R{点}{てん}を\R{読}{よ}む。
\end{enumerate}

\noindent\begin{minipage}[c]{0.52\linewidth}\centering
\begin{tikzpicture}
\begin{axis}[bkplot, width=66mm, height=56mm, xmin=0, xmax=8.6, ymin=0, ymax=34, xtick={0,...,8}, ytick={0,10,20,30},
  xlabel={$t$（s）}, ylabel={\R{位置}{い|ち}（m）}]
  \addplot[bkcurve, domain=0:8] {4*x};
  \addplot[bkcurve2, domain=0:8] {10+2*x};
  \addplot[only marks, mark=*, mark size=1.2pt, bkBlue] coordinates {(0,0) (1,4) (2,8) (3,12) (4,16) (6,24) (7,28)};
  \addplot[only marks, mark=*, mark size=1.2pt, bkOrange] coordinates {(0,10) (1,12) (2,14) (3,16) (4,18) (6,22) (7,24) (8,26)};
  \addplot[only marks, mark=*, mark size=2pt, bkRed] coordinates {(5,20)};
  \node[font=\sffamily\footnotesize, text=bkBlue!80!black, anchor=east, fill=white, inner sep=1pt] at (axis cs:6.4,30) {\R{兄}{あに} $4t$};
  \node[font=\sffamily\footnotesize, text=bkOrange!80!black, anchor=north, fill=white, inner sep=1pt] at (axis cs:7.1,19.5) {\R{弟}{おとうと} $10+2t$};
  \node[font=\sffamily\footnotesize, text=bkRed!80!black, anchor=south east, fill=white, inner sep=1pt] at (axis cs:4.6,21.5) {5\R{秒後}{びょう|ご}に\qty{20}{m}};
\end{axis}
\end{tikzpicture}
\end{minipage}\hfill
\begin{minipage}[c]{0.45\linewidth}
2\R{本}{ほん}の\R{線}{せん}が\R{交}{まじ}わる\R{点}{てん}は、\R{兄}{あに}の\R{線}{せん}の\R{上}{うえ}にも、\R{弟}{おとうと}の\R{線}{せん}の\R{上}{うえ}にもあります。つまり「\R{同}{おな}じ\R{時刻}{じ|こく}に\R{同}{おな}じ\R{位置}{い|ち}にいる」\R{点}{てん}です。

\R{交}{まじ}わる\R{点}{てん}を\R{読}{よ}むと、\R{横}{よこ}が5、たてが20です。\R{表}{ひょう}とも\R{方程式}{ほう|てい|しき}とも\R{同}{おな}じ\R{答}{こた}えになりました。
\end{minipage}

\par\addvspace{6pt}
\textbf{\R{表}{ひょう}では\R{見}{み}のがす\R{答}{こた}え}\quad \R{弟}{おとうと}が\qty{9}{m}\R{先}{さき}からスタートしたら、どうでしょう。\R{表}{ひょう}を\R{作}{つく}ってみます。

\begin{bkcenter}
\begin{bktable}{colspec={Q[l,m,wd=33mm] *{7}{Q[c,m,wd=8mm]}}, rowsep=1pt, colsep=3pt, column{6-7}={bg=bkYellow!20}}
\R{時間}{じ|かん} $t$（s） & 0 & 1 & 2 & 3 & 4 & 5 & 6 \\
\R{兄}{あに} $4t$（m） & 0 & 4 & 8 & 12 & 16 & 20 & 24 \\
\R{弟}{おとうと} $9+2t$（m） & 9 & 11 & 13 & 15 & 17 & 19 & 21 \\
\end{bktable}
\bktabcaption{4\R{秒後}{びょう|ご}はまだ\R{弟}{おとうと}が\R{前}{まえ}（$16<17$）。5\R{秒後}{びょう|ご}は\R{兄}{あに}が\R{前}{まえ}（$20>19$）。}
\end{bkcenter}

\R{表}{ひょう}のどこにも、\R{同}{おな}じ\R{数}{かず}がありません。4\R{秒後}{びょう|ご}と5\R{秒後}{びょう|ご}の\R{間}{あいだ}で\R{追}{お}いこしたことはわかりますが、ちょうど\R{何秒後}{なん|びょう|ご}かはわかりません。\R{方程式}{ほう|てい|しき}なら、ぴったり\R{求}{もと}まります。
\[ 4t=9+2t\ \to\ 4t-2t=9\ \to\ 2t=9\ \to\ t=4.5 \]
\R{答}{こた}えは4.5\R{秒後}{びょう|ご}です（\R{確}{たし}かめ：$4\times4.5=18$、$9+2\times4.5=18$）。

\begin{bkcenter}
\begin{bktable}{colspec={Q[l,m,wd=20mm] Q[l,m,wd=50mm] Q[l,m,wd=48mm]}, rowsep=0.5pt, cells={font=\small}, row{1}={font=\sffamily\bfseries\small}}
\R{方法}{ほう|ほう} & よいところ & こまるところ \\
\R{表}{ひょう} & かけ\R{算}{ざん}とたし\R{算}{ざん}だけでできる & 4.5\R{秒}{びょう}のような\R{答}{こた}えは\R{見}{み}のがす \\
グラフ & ようすが\R{目}{め}で\R{見}{み}える & \R{細}{こま}かい\R{値}{あたい}が\R{読}{よ}みにくい \\
\R{方程式}{ほう|てい|しき} & いつでもぴったり\R{求}{もと}まる & ようすは\R{見}{み}えにくい \\
\end{bktable}
\bktabcaption{3つの\R{方法}{ほう|ほう}をくらべよう。\R{組}{く}み\R{合}{あ}わせて\R{使}{つか}うと\R{強}{つよ}い。}
\end{bkcenter}

\bksubsection{8}{\R{答}{こた}えを\R{確}{たし}かめる}

\leavevmode\label{bk:2-16:tashikame}\R{方程式}{ほう|てい|しき}を\R{解}{と}いたら、\R{解}{かい}をもとの\R{式}{しき}に\R{代入}{だい|にゅう}して、\R{左辺}{さ|へん}と\R{右辺}{う|へん}が\R{等}{ひと}しくなるか\R{確}{たし}かめます。\R{途中}{と|ちゅう}で\R{符号}{ふ|ごう}などをまちがえていたら、ここで\R{気}{き}づけます。

$t=5$ なら、\R{左辺}{さ|へん}は $4\times5=20$、\R{右辺}{う|へん}は $10+2\times5=20$ で\R{等}{ひと}しいので、\R{正}{ただ}しい\R{解}{かい}です。

\R{答}{こた}えが\R{場面}{ば|めん}に\R{合}{あ}うかも\R{見}{み}ます。たとえば、\R{兄}{あに}が\R{秒速}{びょう|そく}\qty{1.5}{m}しか\R{出}{だ}せなかったとします。\R{方程式}{ほう|てい|しき}は $1.5t=10+2t$ です。
\[ 1.5t-2t=10\ \to\ -0.5t=10\ \to\ t=10\div(-0.5)=-20 \]
\R{答}{こた}えはマイナスの20です。\R{時間}{じ|かん}のマイナスは「スタートより\R{前}{まえ}」を\R{表}{あらわ}します（2-2）。もし2\R{人}{にん}がずっと\R{前}{まえ}から\R{走}{はし}っていたとしたら、20\R{秒}{びょう}\R{前}{まえ}に\R{同}{おな}じ\R{場所}{ば|しょ}にいた、という\R{意味}{い|み}です。でも2\R{人}{にん}はスタートから\R{走}{はし}り\R{出}{だ}したので、そのような\R{時}{とき}はありません。\R{兄}{あに}のほうがおそいので、\R{差}{さ}は1\R{秒}{びょう}に\qty{0.5}{m}ずつ\R{広}{ひろ}がり、これから\R{先}{さき}は\R{追}{お}いつけません。\R{式}{しき}の\R{答}{こた}えが\R{出}{で}ても、\R{場面}{ば|めん}にもどして\R{意味}{い|み}を\R{考}{かんが}えることが\R{大切}{たい|せつ}です。

\begin{motto}[2]
\textbf{\R{解}{かい}がない\R{方程式}{ほう|てい|しき}}\par
2\R{人}{にん}とも\R{秒速}{びょう|そく}\qty{2}{m}なら、\R{方程式}{ほう|てい|しき}は $2t=10+2t$ です。$2t$ を\R{移項}{い|こう}すると $2t-2t=10$、つまり $0=10$ になります。$t$ に\R{何}{なに}を\R{入}{い}れても0は10になりません。だから、この\R{方程式}{ほう|てい|しき}には\R{解}{かい}がありません。\R{場面}{ば|めん}にもどすと「2\R{人}{にん}は\R{同}{おな}じ\R{速}{はや}さなので、いつまでも\qty{10}{m}はなれたまま」です。グラフでは、2\R{本}{ほん}の\R{線}{せん}が\R{同}{おな}じかたむきで\R{平行}{へい|こう}になり、いつまでも\R{交}{まじ}わりません（3-6）。
\end{motto}

\begin{taisetsu}
① わからない\R{数}{かず}を\R{文字}{も|じ}にする　② \R{等}{ひと}しい\R{量}{りょう}で\R{方程式}{ほう|てい|しき}をつくる\\
③ \R{移項}{い|こう}・まとめる・わるを1\R{行}{ぎょう}ずつ　④ \R{代入}{だい|にゅう}して\R{確}{たし}かめ、\R{場面}{ば|めん}に\R{合}{あ}うか\R{見}{み}る
\end{taisetsu}

\begin{kotoba}
\small
\bkword{\R{方程式}{ほう|てい|しき}}{\R{文字}{も|じ}に\R{入}{い}れる\R{数}{かず}によって、\R{成}{な}り\R{立}{た}ったり\R{成}{な}り\R{立}{た}たなかったりする\R{等式}{とう|しき}。}
\bkword{\R{解}{かい}}{\R{方程式}{ほう|てい|しき}を\R{成}{な}り\R{立}{た}たせる\R{数}{かず}。\R{解}{かい}を\R{求}{もと}めることを「\R{方程式}{ほう|てい|しき}を\R{解}{と}く」という。}
\bkword{1\R{次方程式}{じ|ほう|てい|しき}}{\R{文字}{も|じ}の\R{項}{こう}が「\R{数}{かず}×\R{文字}{も|じ}」だけの\R{方程式}{ほう|てい|しき}。$ax=b$ の\R{形}{かたち}にして\R{解}{と}ける。}
\end{kotoba}

\clearpage
%% ---------------- 例題 ----------------
\bksection{reidai}{\R{例題}{れい|だい}でたしかめよう}

\begin{reidai}[\R{移項}{い|こう}するだけ・わるだけの\R{方程式}{ほう|てい|しき}]
\R{方程式}{ほう|てい|しき}を\R{解}{と}きましょう。\quad（1）\nobreak $x+7=3$\quad（2）\nobreak $5x=-15$\par
\kangaekata （1）は $+7$ を\R{移項}{い|こう}するだけ、（2）は\R{両辺}{りょう|へん}を5でわるだけです。\R{答}{こた}えがマイナスになることもあります（2-2）。\par
\shiki\begin{minipage}[t]{0.42\linewidth}
（1）\nobreak $x=3-7$\quad $+7$ を\R{移項}{い|こう}\par
\hspace*{1.6\zw}$x=-4$
\end{minipage}\hfill
\begin{minipage}[t]{0.46\linewidth}
（2）\nobreak $x=(-15)\div5$\quad \R{両辺}{りょう|へん}を5でわる\par
\hspace*{1.6\zw}$x=-3$\quad マイナス÷プラスはマイナス
\end{minipage}\par
\kotae （1）\nobreak $x=-4$（\R{確}{たし}かめ：$-4+7=3$）\quad（2）\nobreak $x=-3$（\R{確}{たし}かめ：$5\times(-3)=-15$）
\end{reidai}

\begin{jibun}
\R{方程式}{ほう|てい|しき}を\R{解}{と}き、\R{代入}{だい|にゅう}して\R{確}{たし}かめましょう。\par
（1）\nobreak $x+8=2$\quad（2）\nobreak $4x=-28$
\bkjans{（1）\nobreak $x=2-8=-6$（$-6+8=2$）\quad（2）\nobreak $x=(-28)\div4=-7$（$4\times(-7)=-28$）}
\end{jibun}

\bkfukidashi{mina}{\R{答}{こた}えがマイナスでも、あわてなくていいよ。\R{代入}{だい|にゅう}して\R{左辺}{さ|へん}と\R{右辺}{う|へん}が\R{等}{ひと}しくなれば、それが\R{正}{ただ}しい\R{解}{かい}なんだ。}

\begin{reidai}[\R{移項}{い|こう}して\R{解}{と}く]
\R{方程式}{ほう|てい|しき} $5x-3=2x+9$ を\R{解}{と}きましょう。\par
\shiki\begin{bkiiqhenkei}[t]{}
$5x-3$ & ${}=2x+9$ & \\
$5x-2x$ & ${}=9+3$ & $+2x$ と $-3$ を\R{移項}{い|こう} \\
$3x$ & ${}=12$ & まとめる \\
$x$ & ${}=4$ & \R{両辺}{りょう|へん}を3でわる \\
\end{bkiiqhenkei}\hfill
\begin{minipage}[t]{36mm}
\kangaekata\par
$x$ の\R{項}{こう}は\R{左辺}{さ|へん}へ、\R{数}{かず}の\R{項}{こう}は\R{右辺}{う|へん}へ。\R{移項}{い|こう}すると\R{符号}{ふ|ごう}が\R{変}{か}わる。
\end{minipage}\par
\kotae $x=4$（\R{確}{たし}かめ：\R{左辺}{さ|へん} $5\times4-3=17$、\R{右辺}{う|へん} $2\times4+9=17$）
\end{reidai}

\begin{jibun}
\R{方程式}{ほう|てい|しき} $7x+2=4x+17$ を\R{解}{と}き、\R{代入}{だい|にゅう}して\R{確}{たし}かめましょう。
\bkjans{$7x-4x=17-2$、$3x=15$、$x=5$（\R{両辺}{りょう|へん}とも37）}
\end{jibun}


\begin{reidai}[かっこのある\R{方程式}{ほう|てい|しき}]
\R{方程式}{ほう|てい|しき} $3(x+4)=5x+2$ を\R{解}{と}きましょう。\par
\shiki\begin{bkiiqhenkei}[t]{}
$3(x+4)$ & ${}=5x+2$ & \\
$3x+12$ & ${}=5x+2$ & かっこをはずす \\
$3x-5x$ & ${}=2-12$ & $+5x$ と $+12$ を\R{移項}{い|こう} \\
$-2x$ & ${}=-10$ & まとめる \\
$x$ & ${}=5$ & \R{両辺}{りょう|へん}を $-2$ でわる \\
\end{bkiiqhenkei}\hfill
\begin{minipage}[t]{36mm}
\kangaekata\par
まず\R{分配法則}{ぶん|ぱい|ほう|そく}（2-12）でかっこをはずす。$x$ にも4にも3をかける。$(-10)\div(-2)$ の\R{答}{こた}えはプラス（2-2）。
\end{minipage}\par
\kotae $x=5$（\R{確}{たし}かめ：\R{左辺}{さ|へん} $3\times(5+4)=27$、\R{右辺}{う|へん} $5\times5+2=27$）
\end{reidai}

\begin{reidai}[\R{小数}{しょう|すう}のある\R{方程式}{ほう|てい|しき}]
\R{方程式}{ほう|てい|しき} $0.3x+2=0.1x+3$ を\R{解}{と}きましょう。\par
\shiki\begin{bkiiqhenkei}[t]{}
$0.3x+2$ & ${}=0.1x+3$ & \\
$3x+20$ & ${}=x+30$ & すべての\R{項}{こう}に10をかける \\
$3x-x$ & ${}=30-20$ & \R{移項}{い|こう} \\
$2x$ & ${}=10$ & まとめる \\
$x$ & ${}=5$ & \R{両辺}{りょう|へん}を2でわる \\
\end{bkiiqhenkei}\hfill
\begin{minipage}[t]{36mm}
\kangaekata\par
10\R{倍}{ばい}すると\R{小数点}{しょう|すう|てん}が\R{右}{みぎ}へ1つ\R{動}{うご}く（2-1）。$0.3\times10=3$。\R{等式}{とう|しき}をこわさないよう、2や3にも10をかける。
\end{minipage}\par
\kotae $x=5$（\R{確}{たし}かめ：\R{左辺}{さ|へん} $0.3\times5+2=3.5$、\R{右辺}{う|へん} $0.1\times5+3=3.5$）
\end{reidai}

\begingroup\tcbset{equal height group=bkiiqjB}%
\noindent\begin{minipage}[t]{0.49\linewidth}
\begin{jibun}
\R{方程式}{ほう|てい|しき} $2(x-3)=x+4$ を\R{解}{と}き、\R{代入}{だい|にゅう}して\R{確}{たし}かめましょう。
\bkjans{$2x-6=x+4$、$2x-x=4+6$、$x=10$（\R{両辺}{りょう|へん}とも14）}
\end{jibun}
\end{minipage}\hfill
\begin{minipage}[t]{0.49\linewidth}
\begin{jibun}
\R{方程式}{ほう|てい|しき} $0.2x+1=0.5x-2$ を\R{解}{と}き、\R{代入}{だい|にゅう}して\R{確}{たし}かめましょう。
\bkjans{$2x+10=5x-20$、$2x-5x=-20-10$、$-3x=-30$、$x=10$（\R{両辺}{りょう|へん}とも3）}
\end{jibun}
\end{minipage}\par
\endgroup

\begin{reidai}[\R{分数}{ぶん|すう}のある\R{方程式}{ほう|てい|しき}]
\R{方程式}{ほう|てい|しき} $\dfrac{x}{4}+3=x-6$ を\R{解}{と}きましょう。\par
\kangaekata $\dfrac{x}{4}$ は $x\div4$ です（2-12）。4をかけると $x\div4\times4=x$ にもどり、\R{分数}{ぶん|すう}が\R{消}{き}えます。\R{等式}{とう|しき}をこわさないように、\R{左辺}{さ|へん}にも\R{右辺}{う|へん}にも、\textbf{すべての\R{項}{こう}}に4をかけます（\R{分配法則}{ぶん|ぱい|ほう|そく}）。これを「\R{分母}{ぶん|ぼ}をはらう」といいます。\par
\shiki\begin{bkiiqhenkei}[t]{}
$\dfrac{x}{4}+3$ & ${}=x-6$ & \\
$\dfrac{x}{4}\times4+3\times4$ & ${}=x\times4-6\times4$ & \R{両辺}{りょう|へん}のすべての\R{項}{こう}に4をかける \\
$x+12$ & ${}=4x-24$ & \R{計算}{けい|さん}する \\
$x-4x$ & ${}=-24-12$ & $+4x$ と $+12$ を\R{移項}{い|こう} \\
$-3x$ & ${}=-36$ & まとめる \\
$x$ & ${}=12$ & \R{両辺}{りょう|へん}を $-3$ でわる \\
\end{bkiiqhenkei}\par
\kotae $x=12$（\R{確}{たし}かめ：\R{左辺}{さ|へん} $12\div4+3=6$、\R{右辺}{う|へん} $12-6=6$）
\end{reidai}

\begin{jibun}
\R{方程式}{ほう|てい|しき} $\dfrac{x}{2}+4=x+1$ を\R{解}{と}き、\R{代入}{だい|にゅう}して\R{確}{たし}かめましょう。\R{両辺}{りょう|へん}のすべての\R{項}{こう}に2をかけてから\R{解}{と}こう。
\bkjans{$x+8=2x+2$、$x-2x=2-8$、$-x=-6$、$x=6$（\R{両辺}{りょう|へん}とも7）}
\end{jibun}

\clearpage
\begin{reidai}[\R{追}{お}いかけ（\R{単位}{たん|い}をそろえる）]
\R{分速}{ふん|そく}\qty{60}{m}で\R{歩}{ある}く\R{弟}{おとうと}を、3\R{分後}{ぷん|ご}に\R{兄}{あに}が\R{秒速}{びょう|そく}\qty{2.5}{m}の\R{自転車}{じ|てん|しゃ}で\R{追}{お}いかけました。\R{兄}{あに}は\R{出}{で}てから\R{何分後}{なん|ぷん|ご}に\R{追}{お}いつきますか。\par
\kangaekata ① \R{速}{はや}さの\R{単位}{たん|い}がちがうので、\R{兄}{あに}の\R{速}{はや}さを\R{分速}{ふん|そく}にそろえます。1\R{分}{ぷん}は60\R{秒}{びょう}なので、\qty{60}{s/min}をかけます（2-10）。
\[ \qty{2.5}{m/s}\times\qty{60}{s/min}=2.5\,\frac{\mathrm{m}}{\cancel{\mathrm{s}}}\times\frac{60\,\cancel{\mathrm{s}}}{1\,\mathrm{min}}=\qty{150}{m/min} \]
\hspace*{2.6\zw}sが\R{約分}{やく|ぶん}されて、\R{分速}{ふん|そく}（m/min）になりました。\par
\hspace*{2.6\zw}② \R{兄}{あに}が\R{出}{で}てから $x$ \R{分後}{ふん|ご}に\R{追}{お}いつくとします。\R{弟}{おとうと}は\R{兄}{あに}より3\R{分}{ぷん}\R{早}{はや}く\R{出}{で}たので、そのとき\R{弟}{おとうと}は $x+3$ \R{分}{ふん}\R{歩}{ある}いています（\R{下}{した}の\R{図}{ず}）。\par
\begin{bkfigure}
\begin{tikzpicture}[x=1mm, y=1mm]
  \draw[bkGray!60, line width=0.6pt, -{Stealth[length=4pt]}] (0,-6) -- (110,-6);
  \node[anchor=north east, font=\sffamily\footnotesize, text=bkInk] at (110,-7.5) {\R{時間}{じ|かん}};
  \fill[bkOrange!30] (0,8) rectangle (100,14); \draw[bkOrange!80!black, line width=0.6pt] (0,8) rectangle (100,14);
  \node[font=\sffamily\footnotesize, text=bkInk] at (50,11) {\R{弟}{おとうと}が\R{歩}{ある}く\R{時間}{じ|かん}\enspace $x+3$ \R{分}{ふん}};
  \draw[bkGray, line width=0.6pt, dashed] (0,0) rectangle (40,6);
  \node[font=\sffamily\footnotesize, text=bkGray] at (20,3) {\R{兄}{あに}はまだ\R{家}{いえ}\enspace 3\R{分}{ぷん}};
  \fill[bkBlue!25] (40,0) rectangle (100,6); \draw[bkBlue!80!black, line width=0.6pt] (40,0) rectangle (100,6);
  \node[font=\sffamily\footnotesize, text=bkInk] at (70,3) {\R{兄}{あに}が\R{走}{はし}る\R{時間}{じ|かん}\enspace $x$ \R{分}{ふん}};
  \foreach \p/\t in {0/{\R{弟}{おとうと}が\R{出}{で}る}, 40/{\R{兄}{あに}が\R{出}{で}る}, 100/{\R{追}{お}いつく}} {
    \draw[bkInk, line width=0.6pt] (\p,-7.5) -- (\p,-4.5);
    \node[anchor=south, font=\sffamily\footnotesize, text=bkInk] at (\p,15.5) {\t};}
\end{tikzpicture}
\end{bkfigure}
\hspace*{2.6\zw}③ \R{追}{お}いつくときは、2\R{人}{にん}が\R{家}{いえ}から\R{進}{すす}んだ\R{道}{みち}のりが\R{等}{ひと}しい。\R{道}{みち}のり＝\R{速}{はや}さ×\R{時間}{じ|かん}です。\par
\shiki\begin{bkiiqhenkei}[t]{}
$150x$ & ${}=60(x+3)$ & \R{兄}{あに}の\R{道}{みち}のり＝\R{弟}{おとうと}の\R{道}{みち}のり \\
$150x$ & ${}=60x+180$ & かっこをはずす \\
$150x-60x$ & ${}=180$ & $+60x$ を\R{移項}{い|こう} \\
$90x$ & ${}=180$ & まとめる \\
$x$ & ${}=2$ & \R{両辺}{りょう|へん}を90でわる \\
\end{bkiiqhenkei}\par
\kotae 2\R{分後}{ふん|ご}（\R{確}{たし}かめ：\R{兄}{あに} $150\times2=300$、\R{弟}{おとうと} $60\times(2+3)=300$ で、どちらも\qty{300}{m}）
\end{reidai}

\begin{jibun}
\R{普通}{ふ|つう}\R{列車}{れっ|しゃ}（\R{分速}{ふん|そく}\qty{1}{km}）が\R{駅}{えき}を\R{出}{で}た4\R{分後}{ぷん|ご}に、\R{特急}{とっ|きゅう}（\R{分速}{ふん|そく}\qty{1.5}{km}）が\R{同}{おな}じ\R{駅}{えき}を\R{出}{で}ました。\R{特急}{とっ|きゅう}が\R{出}{で}てから\R{何分後}{なん|ぷん|ご}に\R{追}{お}いつきますか。\R{時間}{じ|かん}の\R{図}{ず}をノートにかいてから\R{式}{しき}をつくろう。
\bkjans{$1.5x=1\times(x+4)$、$1.5x-x=4$、$0.5x=4$、$x=8$ で8\nobreak\mbox{\R{分後}{ぷん|ご}}（\R{駅}{えき}から\qty{12}{km}）}
\end{jibun}

\begin{reidai}[2つの\R{水}{みず}そう]
\R{水}{みず}そうAは\R{深}{ふか}さ\qty{10}{cm}から\R{毎分}{まい|ふん}\qty{2}{cm}ずつ\R{深}{ふか}くなり、Bは\qty{40}{cm}から\R{毎分}{まい|ふん}\qty{1}{cm}ずつ\R{浅}{あさ}くなります。\R{深}{ふか}さが\R{同}{おな}じになるのは\R{何分後}{なん|ぷん|ご}で、\R{何}{なん}cmですか。\par
\kangaekata $x$ \R{分後}{ふん|ご}の\R{深}{ふか}さを\R{式}{しき}にします。Aは「はじめの10」に「ふえた $2x$」をたして $10+2x$。Bは「はじめの40」から「へった $1\times x=x$」をひいて $40-x$。この2つが\R{等}{ひと}しいとおきます。\par
\shiki\begin{bkiiqhenkei}[t]{}
$10+2x$ & ${}=40-x$ & \\
$2x+x$ & ${}=40-10$ & $-x$ と $+10$ を\R{移項}{い|こう} \\
$3x$ & ${}=30$ & まとめる \\
$x$ & ${}=10$ & \R{両辺}{りょう|へん}を3でわる \\
\end{bkiiqhenkei}\par
\kotae 10\R{分後}{ぷん|ご}、\qty{30}{cm}（\R{確}{たし}かめ：A $10+2\times10=30$、B $40-10=30$）
\end{reidai}

\begin{jibun}
Aは\qty{5}{cm}から\R{毎分}{まい|ふん}\qty{3}{cm}、Bは\qty{25}{cm}から\R{毎分}{まい|ふん}\qty{1}{cm}ずつ\R{深}{ふか}くなります。\R{深}{ふか}さが\R{同}{おな}じになるのは\R{何分後}{なん|ぷん|ご}で、\R{何}{なん}cmですか。
\bkjans{$5+3x=25+x$、$2x=20$、$x=10$ で10\nobreak\mbox{\R{分後}{ぷん|ご}}、\qty{35}{cm}}
\end{jibun}

\begin{reidai}[あまる・たりない]
えんぴつを\R{子}{こ}どもに\R{配}{くば}ります。\R{1人}{ひとり}2\R{本}{ほん}ずつ\R{配}{くば}ると5\R{本}{ほん}あまり、3\R{本}{ぼん}ずつ\R{配}{くば}ると4\R{本}{ほん}たりません。\R{子}{こ}どもの\R{人数}{にん|ずう}と、えんぴつの\R{本数}{ほん|すう}を\R{求}{もと}めましょう。\par
\kangaekata \R{子}{こ}どもを $x$ \R{人}{にん}とします。えんぴつの\R{本数}{ほん|すう}を、2\R{通}{とお}りの\R{配}{くば}り\R{方}{かた}で\R{表}{あらわ}します。\par
\hspace*{2.6\zw}・2\R{本}{ほん}ずつ：\R{配}{くば}った $2x$ \R{本}{ほん}に、あまりの5\R{本}{ほん}をたすと\R{全部}{ぜん|ぶ}。$2x+5$\par
\hspace*{2.6\zw}・3\R{本}{ぼん}ずつ：\R{配}{くば}るのに $3x$ \R{本}{ほん}いるのに、4\R{本}{ほん}たりない。\R{全部}{ぜん|ぶ}は $3x$ より4\R{少}{すく}ない $3x-4$\par
\begin{bkfigure}
\begin{tikzpicture}[x=1mm, y=1mm]
  \fill[bkBlue!22] (0,10) rectangle (70,16); \draw[bkInk, line width=0.6pt] (0,10) rectangle (70,16);
  \node[font=\sffamily\footnotesize] at (35,13) {2\R{本}{ほん}ずつ $x$ \R{人分}{にん|ぶん}\enspace $2x$};
  \fill[bkGreen!25] (70,10) rectangle (88,16); \draw[bkInk, line width=0.6pt] (70,10) rectangle (88,16);
  \node[font=\sffamily\footnotesize] at (79,13) {あまり5};
  \fill[bkOrange!22] (0,0) rectangle (88,6); \draw[bkInk, line width=0.6pt] (0,0) rectangle (88,6);
  \draw[bkInk, line width=0.6pt, dashed] (88,0) rectangle (106,6);
  \node[font=\sffamily\footnotesize] at (44,3) {\R{持}{も}っているえんぴつ（$3x$ より4\R{少}{すく}ない）};
  \node[font=\sffamily\footnotesize, text=bkRed!80!black] at (97,3) {たりない4};
  \draw[bkRed, line width=0.8pt] (88,-2) -- (88,18);
  \draw[decorate, decoration={brace, amplitude=2mm}, line width=0.6pt] (0,18.5) -- (88,18.5);
  \node[anchor=south, font=\sffamily\footnotesize] at (44,21.5) {えんぴつ\R{全部}{ぜん|ぶ}（どちらも\R{同}{おな}じ\R{本数}{ほん|すう}）};
  \draw[decorate, decoration={brace, amplitude=2mm, mirror}, line width=0.6pt] (0,-2) -- (106,-2);
  \node[anchor=north, font=\sffamily\footnotesize] at (53,-5) {3\R{本}{ぼん}ずつ\R{配}{くば}るのに\R{必要}{ひつ|よう}な\R{本数}{ほん|すう} $3x$};
\end{tikzpicture}
\end{bkfigure}
\shiki\begin{bkiiqhenkei}[t]{}
$2x+5$ & ${}=3x-4$ & どちらもえんぴつ\R{全部}{ぜん|ぶ}の\R{本数}{ほん|すう} \\
$2x-3x$ & ${}=-4-5$ & $+3x$ と $+5$ を\R{移項}{い|こう} \\
$-x$ & ${}=-9$ & まとめる \\
$x$ & ${}=9$ & \R{両辺}{りょう|へん}を $-1$ でわる \\
\end{bkiiqhenkei}\par
\kotae 9\R{人}{にん}、えんぴつは $2\times9+5=23$ で23\R{本}{ぼん}（\R{確}{たし}かめ：$3\times9-4=23$）
\end{reidai}

\begin{jibun}
\R{折}{お}り\R{紙}{がみ}を\R{1人}{ひとり}4\R{枚}{まい}ずつ\R{配}{くば}ると2\R{枚}{まい}あまり、6\R{枚}{まい}ずつ\R{配}{くば}ると10\R{枚}{まい}たりません。\R{人数}{にん|ずう}と\R{折}{お}り\R{紙}{がみ}の\R{枚数}{まい|すう}を\R{求}{もと}めましょう。
\bkjans{$4x+2=6x-10$、$4x-6x=-10-2$、$-2x=-12$、$x=6$。6\nobreak\R{人}{にん}、$4\times6+2=26$ で26\nobreak\R{枚}{まい}}
\end{jibun}

\begin{reidai}[グラフで\R{解}{と}く]
$y=2t$ と $y=6+t$ のグラフをかき、\R{交}{まじ}わる\R{点}{てん}を\R{読}{よ}みとりましょう。また、$2t=6+t$ を\R{解}{と}いて、\R{同}{おな}じ\R{答}{こた}えになるか\R{確}{たし}かめましょう。\par
\kangaekata ① $t=0$、1、2、…を\R{代入}{だい|にゅう}して\R{表}{ひょう}を\R{作}{つく}る　② \R{点}{てん}をとる　③ \R{線}{せん}でつなぐ　④ \R{交}{まじ}わる\R{点}{てん}を\R{読}{よ}む。\par
\noindent\begin{minipage}[c]{0.5\linewidth}\centering
\begin{bktable}{colspec={Q[l,m,wd=14mm] *{5}{Q[c,m,wd=6mm]}}, rowsep=0.5pt, colsep=2pt, cells={font=\small}}
$t$ & 0 & 2 & 4 & 6 & 8 \\
$2t$ & 0 & 4 & 8 & 12 & 16 \\
$6+t$ & 6 & 8 & 10 & 12 & 14 \\
\end{bktable}
\par\vspace{4pt}
\begin{bkiiqhenkei}{}
$2t$ & ${}=6+t$ & \\
$2t-t$ & ${}=6$ & $+t$ を\R{移項}{い|こう} \\
$t$ & ${}=6$ & まとめる \\
\end{bkiiqhenkei}
\end{minipage}\hfill
\begin{minipage}[c]{0.47\linewidth}\centering
\begin{tikzpicture}
\begin{axis}[bkplot, width=56mm, height=46mm, xmin=0, xmax=9.4, ymin=0, ymax=19, xtick={0,2,...,8}, ytick={0,4,...,16},
  xlabel={$t$}, ylabel={$y$}]
  \addplot[bkcurve, domain=0:9] {2*x};
  \addplot[bkcurve2, domain=0:9] {6+x};
  \addplot[only marks, mark=*, mark size=1.2pt, bkBlue] coordinates {(0,0) (2,4) (4,8) (8,16)};
  \addplot[only marks, mark=*, mark size=1.2pt, bkOrange] coordinates {(0,6) (2,8) (4,10) (8,14)};
  \addplot[only marks, mark=*, mark size=2pt, bkRed] coordinates {(6,12)};
  \node[font=\sffamily\footnotesize, text=bkBlue!80!black, anchor=east, fill=white, inner sep=1pt] at (axis cs:7.4,17) {$y=2t$};
  \node[font=\sffamily\footnotesize, text=bkOrange!80!black, anchor=north, fill=white, inner sep=1pt] at (axis cs:8.2,12.5) {$y=6+t$};
\end{axis}
\end{tikzpicture}
\end{minipage}\par
\kotae \R{交}{まじ}わる\R{点}{てん}は $t=6$、$y=12$。\R{方程式}{ほう|てい|しき}の\R{解}{かい}も $t=6$ で\R{同}{おな}じ
\end{reidai}

\begin{jibun}
$y=4t$ と $y=9+t$ のグラフをノートにかき、\R{交}{まじ}わる\R{点}{てん}を\R{読}{よ}みとりましょう。また、$4t=9+t$ を\R{解}{と}いて、\R{同}{おな}じ\R{答}{こた}えになるか\R{確}{たし}かめましょう。
\bkjans{\R{交}{まじ}わるのは $t=3$、$y=12$。$4t-t=9$、$3t=9$、$t=3$ で\R{同}{おな}じ}
\end{jibun}

\bkfukidashi{delta}{どの\R{例題}{れい|だい}も、「わからない\R{数}{かず}を\R{文字}{も|じ}にする→\R{等}{ひと}しい\R{量}{りょう}で\R{式}{しき}をつくる→\R{移項}{い|こう}・まとめる・わる→\R{確}{たし}かめる」の\R{順}{じゅん}だったね。かっこ・\R{分数}{ぶん|すう}・\R{小数}{しょう|すう}は、はじめに\R{形}{かたち}を\R{整}{ととの}えよう。}

\begin{reidai}[\R{向}{む}かい\R{合}{あ}って\R{出会}{で|あ}う]
\R{家}{いえ}と\R{公園}{こう|えん}は\qty{900}{m}はなれています。\R{姉}{あね}は\R{家}{いえ}から\R{分速}{ふん|そく}\qty{70}{m}で、\R{弟}{おとうと}は\R{公園}{こう|えん}から\R{分速}{ふん|そく}\qty{80}{m}で、\R{同}{おな}じときに\R{向}{む}かい\R{合}{あ}って\R{歩}{ある}き\R{出}{だ}しました。2\R{人}{にん}は\R{何分後}{なん|ぷん|ご}に\R{出会}{で|あ}いますか。\par
\begin{bkfigure}
\begin{tikzpicture}[x=1mm, y=1mm]
  \draw[bkGray!60, line width=3pt] (0,0) -- (108,0);
  \draw[bkInk, line width=0.7pt] (0,-1.6) -- (0,1.6); \draw[bkInk, line width=0.7pt] (108,-1.6) -- (108,1.6);
  \node[anchor=north, font=\sffamily\footnotesize] at (0,-2.5) {\R{家}{いえ}};
  \node[anchor=north, font=\sffamily\footnotesize] at (108,-2.5) {\R{公園}{こう|えん}};
  \fill[bkBlue!25] (0,4) rectangle (50.4,9); \draw[bkBlue!80!black, line width=0.6pt] (0,4) rectangle (50.4,9);
  \node[font=\sffamily\footnotesize] at (25.2,6.5) {\R{姉}{あね} $70x$ m};
  \fill[bkOrange!25] (50.4,4) rectangle (108,9); \draw[bkOrange!80!black, line width=0.6pt] (50.4,4) rectangle (108,9);
  \node[font=\sffamily\footnotesize] at (79.2,6.5) {\R{弟}{おとうと} $80x$ m};
  \draw[bk arrow, bkBlue] (2,13) -- (20,13);
  \draw[bk arrow, bkOrange] (106,13) -- (88,13);
  \draw[{Stealth[length=4pt]}-{Stealth[length=4pt]}, bkInk, line width=0.6pt] (0,-9.5) -- (108,-9.5);
  \node[anchor=north, font=\sffamily\footnotesize] at (54,-10.5) {\qty{900}{m}};
\end{tikzpicture}
\end{bkfigure}
\kangaekata $x$ \R{分後}{ふん|ご}に\R{出会}{で|あ}うとします。そのとき、\R{姉}{あね}が\R{歩}{ある}いた $70x$ と\R{弟}{おとうと}が\R{歩}{ある}いた $80x$ をつなぐと、ちょうど\qty{900}{m}になります（\R{図}{ず}）。\R{追}{お}いかけとちがい、\R{向}{む}かい\R{合}{あ}うときは2\R{人}{にん}とも\R{間}{あいだ}をちぢめるので、\R{近}{ちか}づく\R{速}{はや}さは\R{和}{わ}の $70+80=150$（m/min）です。\par
\shiki\begin{bkiiqhenkei}[t]{}
$70x+80x$ & ${}=900$ & 2\R{人}{にん}の\R{道}{みち}のりの\R{和}{わ}が\qty{900}{m} \\
$150x$ & ${}=900$ & まとめる \\
$x$ & ${}=6$ & \R{両辺}{りょう|へん}を150でわる \\
\end{bkiiqhenkei}\par
\kotae 6\R{分後}{ふん|ご}（\R{確}{たし}かめ：$70\times6+80\times6=420+480=900$。\R{差}{さ}の\R{考}{かんが}え\R{方}{かた}でも $\qty{900}{m}\div\qty{150}{m/min}=\qty{6}{min}$）
\end{reidai}

\begin{jibun}
\qty{1200}{m}はなれた2つの\R{場所}{ば|しょ}から、\R{妹}{いもうと}は\R{分速}{ふん|そく}\qty{60}{m}、\R{兄}{あに}は\R{分速}{ふん|そく}\qty{100}{m}で、\R{同}{おな}じときに\R{向}{む}かい\R{合}{あ}って\R{歩}{ある}き\R{出}{だ}しました。2\R{人}{にん}は\R{何分後}{なん|ぷん|ご}に\R{出会}{で|あ}いますか。
\bkjans{$60x+100x=1200$、$160x=1200$、$x=7.5$ で7.5\nobreak\mbox{\R{分後}{ふん|ご}}}
\end{jibun}


\clearpage
%% ---------------- まちがえやすいところ ----------------
\bksection{tsumazuki}{まちがえやすいところをチェック}

\begin{tsumazuki}{\R{移項}{い|こう}で\R{符号}{ふ|ごう}を\R{変}{か}えない}
\bkng{$4t=10+2t$ から\quad $4t+2t=10$}
\bkok{$4t=10+2t$ から\quad $4t-2t=10$}
\R{移項}{い|こう}は「\R{両辺}{りょう|へん}から $2t$ をひく」ことです。まちがえたまま $6t=10$ で\R{解}{と}くと $t=\dfrac{5}{3}$ になり、\R{代入}{だい|にゅう}すると\R{左辺}{さ|へん} $4\times\dfrac{5}{3}=\dfrac{20}{3}$、\R{右辺}{う|へん} $10+2\times\dfrac{5}{3}=\dfrac{40}{3}$ で\R{等}{ひと}しくなりません。\R{確}{たし}かめをすれば、まちがいに\R{気}{き}づけます。
\end{tsumazuki}

\begin{tsumazuki}{マイナスでわるとき、\R{符号}{ふ|ごう}をまちがえる}
\bkng{$-2x=-10$ から\quad $x=-5$}
\bkok{$-2x=-10$ から\quad $x=(-10)\div(-2)=5$}
マイナスどうしのわり\R{算}{ざん}の\R{答}{こた}えはプラスです（2-2）。\R{係数}{けい|すう}の $-2$ でわるときは、\R{符号}{ふ|ごう}もいっしょにわります。
\end{tsumazuki}

\begin{tsumazuki}{かっこの\R{後}{うし}ろの\R{項}{こう}にかけわすれる}
\bkng{$3(x+4)=3x+4$}
\bkok{$3(x+4)=3x+12$}
かっこの\R{中}{なか}のどの\R{項}{こう}にも3をかけます（2-12）。たて3、\R{横}{よこ} $x+4$ の\R{長方形}{ちょう|ほう|けい}の\R{面積}{めん|せき}で\R{考}{かんが}えると、\R{横}{よこ}の $x$ の\R{部分}{ぶ|ぶん}も4の\R{部分}{ぶ|ぶん}もたてが3なので、まちがえません。
\end{tsumazuki}

\begin{tsumazuki}{\R{分母}{ぶん|ぼ}をはらうとき、\R{一部}{いち|ぶ}の\R{項}{こう}にしかかけない}
\bkng{$\dfrac{x}{4}+3=x-6$ から\quad $x+3=4x-24$}
\bkok{$\dfrac{x}{4}+3=x-6$ から\quad $x+12=4x-24$}
\R{両辺}{りょう|へん}に4をかけるときは、\R{分数}{ぶん|すう}でない\R{項}{こう}（3）にも4をかけます。\R{天}{てん}びんの\R{両方}{りょう|ほう}の\R{皿}{さら}の\R{物}{もの}を、\textbf{\R{全部}{ぜん|ぶ}}4\R{倍}{ばい}にするのと\R{同}{おな}じです。
\end{tsumazuki}

\begin{tsumazuki}{「たりない」をたしてしまう}
\bkng{3\R{本}{ぼん}ずつ\R{配}{くば}ると4\R{本}{ほん}たりない\quad → えんぴつは $3x+4$}
\bkok{3\R{本}{ぼん}ずつ\R{配}{くば}ると4\R{本}{ほん}たりない\quad → えんぴつは $3x-4$}
たりないのは、\R{配}{くば}るのに\R{必要}{ひつ|よう}な $3x$ \R{本}{ほん}より、\R{持}{も}っている\R{本数}{ほん|すう}が\R{少}{すく}ないということです。テープの\R{図}{ず}をかくと、まちがえません。
\end{tsumazuki}

\clearpage
%% ---------------- 物理・工学 ----------------
\begin{butsuri}
\textbf{\R{自動}{じ|どう}ブレーキ}\par
\R{車}{くるま}の\R{自動}{じ|どう}ブレーキは、カメラやレーダーで\R{前}{まえ}の\R{車}{くるま}までの\R{距離}{きょ|り}と\R{速}{はや}さをはかり、「このままだと、あと\R{何秒}{なん|びょう}で\R{追}{お}いつくか（ぶつかるか）」を\R{計算}{けい|さん}しています。

\R{自分}{じ|ぶん}の\R{車}{くるま}が\R{時速}{じ|そく}\qty{72}{km}、\qty{30}{m}\R{先}{さき}の\R{車}{くるま}が\R{時速}{じ|そく}\qty{54}{km}で\R{走}{はし}っているとします。\R{距離}{きょ|り}はmなので、\R{速}{はや}さも\R{秒速}{びょう|そく}（m/s）にそろえます。1 kmは1000 mなので\qty{1000}{m/km}、1\R{時間}{じ|かん}は3600\R{秒}{びょう}なので\qty{3600}{s/h}を\R{使}{つか}います（2-10）。
\[ \qty{72}{km/h}\times\qty{1000}{m/km}\div\qty{3600}{s/h}=\frac{72\times1000}{3600}\,\mathrm{m/s}=\qty{20}{m/s} \]
\[ \qty{54}{km/h}\times\qty{1000}{m/km}\div\qty{3600}{s/h}=\frac{54\times1000}{3600}\,\mathrm{m/s}=\qty{15}{m/s} \]
kmどうし、hどうしが\R{約分}{やく|ぶん}されて、m/sが\R{残}{のこ}ります。$t$ \R{秒後}{びょう|ご}に\R{追}{お}いつくとすると、
\[ 20t=30+15t\quad\to\quad 20t-15t=30\quad\to\quad 5t=30\quad\to\quad t=6 \]
このままだと6\R{秒}{びょう}でぶつかるので、\R{車}{くるま}は\R{早}{はや}めに\R{音}{おと}で\R{知}{し}らせたり、ブレーキをかけたりします。\R{本物}{ほん|もの}の\R{車}{くるま}は、この\R{計算}{けい|さん}を1\R{秒}{びょう}に\R{何十回}{なん|じっ|かい}もくり\R{返}{かえ}しています。

\begin{bkfigure}
\begin{tikzpicture}[x=1mm, y=1mm]
  \fill[bkGray!25] (-4,-2) rectangle (120,0);
  \bkcar{(0,0)}{0.85}{bkBlue}
  \bkcar{(78,0)}{0.85}{bkRed!80}
  \draw[bk arrow, bkBlue] (28,7) -- (44,7);
  \node[font=\sffamily\footnotesize, text=bkBlue!80!black, anchor=south] at (36,8.5) {\qty{20}{m/s}};
  \draw[bk arrow, bkRed] (105,7) -- (115,7);
  \node[font=\sffamily\footnotesize, text=bkRed!80!black, anchor=south] at (110,8.5) {\qty{15}{m/s}};
  \draw[{Stealth[length=4pt]}-{Stealth[length=4pt]}, bkInk, line width=0.7pt] (25,-5.5) -- (75,-5.5);
  \node[font=\sffamily\footnotesize, text=bkInk, anchor=north] at (50,-6.5) {\qty{30}{m}};
  \foreach \a in {-12,0,12} \draw[bkOrange, line width=0.6pt] (25.5,5) -- ++(\a:5);
\end{tikzpicture}
\bkfigcaption{\R{距離}{きょ|り}と\R{速}{はや}さから、\R{追}{お}いつくまでの\R{時間}{じ|かん}を\R{計算}{けい|さん}する。}
\end{bkfigure}

\end{butsuri}

\begin{butsuri}
\noindent\begin{minipage}[c]{0.53\linewidth}
\textbf{\R{電車}{でん|しゃ}のダイヤ}\par
\R{鉄道}{てつ|どう}の\R{会社}{かい|しゃ}は、\R{横}{よこ}に\R{時刻}{じ|こく}、たてに\R{駅}{えき}の\R{位置}{い|ち}をとったグラフ（ダイヤグラム）で\R{運行}{うん|こう}を\R{計画}{けい|かく}します。\R{線}{せん}が\R{交}{まじ}わる\R{所}{ところ}が、\R{特急}{とっ|きゅう}が\R{普通}{ふ|つう}\R{列車}{れっ|しゃ}に\R{追}{お}いつく\R{所}{ところ}です。\R{線路}{せん|ろ}が1\R{本}{ぽん}の\R{所}{ところ}で\R{追}{お}いついてはこまるので、その\R{前}{まえ}に\R{普通}{ふ|つう}\R{列車}{れっ|しゃ}をB\R{駅}{えき}で\R{待}{ま}たせて、\R{特急}{とっ|きゅう}に\R{先}{さき}へ\R{行}{い}ってもらいます（\R{横}{よこ}に\R{平}{たい}らな\R{線}{せん}が「\R{止}{と}まって\R{待}{ま}つ」）。
\end{minipage}\hfill
\begin{minipage}[c]{0.45\linewidth}\centering
\begin{tikzpicture}[x=1mm, y=1mm]
  \foreach \y/\s in {0/A,12/B,24/C,36/D} {\draw[bkLine, line width=0.4pt] (0,\y) -- (44,\y);
    \node[font=\sffamily\scriptsize, text=bkInk, anchor=east] at (-1.5,\y) {\s\R{駅}{えき}};}
  \draw[bkOrange, line width=1.4pt] (0,0) -- (11,12) -- (20,12) -- (31,24) -- (42,36);
  \draw[bkBlue, line width=1.4pt] (8,0) -- (20,24) -- (26,36);
  \draw[bk axis] (0,0) -- (47,0);
  \node[font=\sffamily\scriptsize, text=bkInk, anchor=north] at (42,-1.5) {\R{時刻}{じ|こく}};
  \node[font=\sffamily\scriptsize, text=bkOrange!80!black, anchor=north west] at (33,22) {\R{普通}{ふ|つう}};
  \node[font=\sffamily\scriptsize, text=bkBlue!80!black, anchor=east] at (20,31) {\R{特急}{とっ|きゅう}};
\end{tikzpicture}
\end{minipage}

\par\addvspace{6pt}
\textbf{\R{大学}{だい|がく}への\R{種}{たね}}\par
\R{表}{ひょう}でためして\R{答}{こた}えに\R{近}{ちか}づく\R{方法}{ほう|ほう}は、コンピュータで\R{速}{はや}くすると、\R{手}{て}では\R{解}{と}けない\R{方程式}{ほう|てい|しき}も\R{解}{と}ける\R{強}{つよ}い\R{道具}{どう|ぐ}になります（2-19、6-2）。\R{大学}{だい|がく}では「\R{数値計算}{すう|ち|けい|さん}」として\R{学}{まな}びます。また、\R{電気}{でん|き}の\R{回路}{かい|ろ}では、わからない\R{電流}{でん|りゅう}をいくつも\R{文字}{も|じ}にして、\R{方程式}{ほう|てい|しき}を\R{何本}{なん|ぼん}もいっしょに\R{解}{と}きます（\R{連立方程式}{れん|りつ|ほう|てい|しき}、3-6）。
\end{butsuri}

\begin{column}{「\R{方程式}{ほう|てい|しき}」ということば}
「\R{方程式}{ほう|てい|しき}」は、2000\R{年}{ねん}ほど\R{前}{まえ}の\R{中国}{ちゅう|ごく}の\R{算数}{さん|すう}の\R{本}{ほん}『\R{九章算術}{きゅう|しょう|さん|じゅつ}』にある「\R{方程}{ほう|てい}」という\R{章}{しょう}の\R{名前}{な|まえ}から\R{来}{き}ています。この\R{本}{ほん}では、\R{数}{かず}を\R{四角}{し|かく}くならべて\R{計算}{けい|さん}していました。「\R{方}{ほう}」は\R{四角}{し|かく}くならべること、「\R{程}{てい}」はくらべてはかることを\R{表}{あらわ}す、といわれています。むかしの\R{人}{ひと}も、わからない\R{数}{かず}を\R{求}{もと}めるために、\R{同}{おな}じように\R{工夫}{く|ふう}していたのです。
\end{column}

\clearpage
%% ---------------- 章末問題 ----------------
\bkendproblems
\bkproblemgroup{\R{基本}{き|ほん}}{2}

\bkproblem \R{次}{つぎ}の\R{方程式}{ほう|てい|しき}を\R{解}{と}きましょう。\par
\hspace*{1\zw}（1）\nobreak $x+9=4$\quad（2）\nobreak $6x=-18$\quad（3）\nobreak $3x-4=x+8$\par
\hspace*{1\zw}（4）\nobreak $2(x+5)=4x$\quad（5）\nobreak $\dfrac{x}{3}+2=x-4$\quad（6）\nobreak $0.4x-1=0.1x+2$
\bkpans{（1）\nobreak $x=-5$\quad（2）\nobreak $x=-3$\quad（3）\nobreak $x=6$\quad（4）\nobreak $x=5$\quad（5）\nobreak $x=9$\quad（6）\nobreak $x=10$}

\bkproblem 2、3、4のうち、\R{方程式}{ほう|てい|しき} $5x-4=2x+5$ の\R{解}{かい}はどれですか。\R{代入}{だい|にゅう}して\R{確}{たし}かめましょう。
\bkpans{3（\R{左辺}{さ|へん}も\R{右辺}{う|へん}も11）}

\bkproblem \R{妹}{いもうと}が\R{家}{いえ}を\R{出}{で}て5\R{分後}{ふん|ご}に、\R{姉}{あね}が\R{自転車}{じ|てん|しゃ}で\R{追}{お}いかけました。\R{妹}{いもうと}は\R{分速}{ふん|そく}\qty{70}{m}、\R{姉}{あね}は\R{秒速}{びょう|そく}\qty{3.5}{m}です。\R{姉}{あね}は\R{出}{で}てから\R{何分後}{なん|ぷん|ご}に\R{追}{お}いつきますか。
\bkpans{$\qty{3.5}{m/s}\times\qty{60}{s/min}=\qty{210}{m/min}$。$210x=70(x+5)$、$140x=350$、$x=2.5$ で2.5\nobreak\mbox{\R{分後}{ふん|ご}}}

\bkproblem \R{水}{みず}そうAには\qty{20}{L}の\R{水}{みず}があり、\R{毎分}{まい|ふん}\qty{4}{L}ずつ\R{入}{い}れます。Bには\qty{80}{L}あり、\R{毎分}{まい|ふん}\qty{2}{L}ずつ\R{出}{だ}します。\R{水}{みず}の\R{量}{りょう}が\R{同}{おな}じになるのは\R{何分後}{なん|ぷん|ご}で、\R{何}{なん}Lですか。
\bkpans{$20+4x=80-2x$、$6x=60$、$x=10$ で10\nobreak\mbox{\R{分後}{ぷん|ご}}、\qty{60}{L}}

\bkproblem $y=3t$ と $y=12+t$ のグラフをノートにかき、\R{交}{まじ}わる\R{点}{てん}を\R{読}{よ}みとりましょう。また、$3t=12+t$ を\R{解}{と}いて、\R{同}{おな}じ\R{答}{こた}えになるか\R{確}{たし}かめましょう。
\bkpans{\R{交}{まじ}わるのは $t=6$ のとき（$y=18$）。$2t=12$ より $t=6$ で\R{同}{おな}じ}
\begin{bkfigure}
\begin{tikzpicture}[x=8mm, y=1.2mm]
  \foreach \x in {1,...,9} \draw[bk grid] (\x,0) -- (\x,31);
  \foreach \y in {5,10,...,30} \draw[bk grid] (0,\y) -- (9.4,\y);
  \draw[bk axis] (0,0) -- (9.8,0);
  \draw[bk axis] (0,0) -- (0,34);
  \foreach \x in {1,...,9} \node[bk tick, below=1.2mm] at (\x,0) {\x};
  \foreach \y in {10,20,30} \node[bk tick, left=1.2mm] at (0,\y) {\y};
  \node[bk tick, below left=1mm] at (0,0) {0};
  \node[anchor=west, bk label] at (9.9,0) {$t$};
  \node[anchor=south, bk label] at (0,34.5) {$y$};
  \fill[bkBlue] (0,0) circle (1.8pt) (1,3) circle (1.8pt) (2,6) circle (1.8pt);
  \fill[bkOrange] (0,12) circle (1.8pt) (1,13) circle (1.8pt) (2,14) circle (1.8pt);
  \node[font=\sffamily\footnotesize, text=bkBlue!80!black, anchor=west, fill=white, inner sep=1pt] at (2.2,5) {$y=3t$};
  \node[font=\sffamily\footnotesize, text=bkOrange!80!black, anchor=west, fill=white, inner sep=1pt] at (2.2,16.5) {$y=12+t$};
\end{tikzpicture}
\bkfigcaption{ノートにかくグラフの\R{例}{れい}。$t=0$、1、2 の\R{点}{てん}をとったところ。\R{続}{つづ}きをかこう。}
\end{bkfigure}

\clearpage
\bkproblemgroup{チャレンジ}{2}

\bkproblem \qty{40}{m}\R{先}{さき}を\qty{15}{m/s}で\R{走}{はし}る\R{車}{くるま}に、\qty{25}{m/s}で\R{近}{ちか}づいています。\R{同}{おな}じ\R{速}{はや}さのままだと、\R{何秒後}{なん|びょう|ご}に\R{追}{お}いつきますか。\R{方程式}{ほう|てい|しき}と、「\R{差}{さ}がちぢまる\R{速}{はや}さ」の2つの\R{方法}{ほう|ほう}で\R{求}{もと}めましょう。
\bkpans{\R{方程式}{ほう|てい|しき}：$25t=40+15t$、$10t=40$、$t=4$。\R{差}{さ}：$\qty{25}{m/s}-\qty{15}{m/s}=\qty{10}{m/s}$、$\qty{40}{m}\div\qty{10}{m/s}=\qty{4}{s}$。どちらも4\nobreak\mbox{\R{秒後}{びょう|ご}}}
\begin{bkfigure}
\begin{tikzpicture}[x=1mm, y=1mm]
  \fill[bkGray!25] (-4,-2) rectangle (110,0);
  \bkcar{(0,0)}{0.7}{bkGreen}
  \bkcar{(70,0)}{0.7}{bkPurple!70}
  \draw[bk arrow, bkGreen] (22,6) -- (36,6);
  \node[font=\sffamily\footnotesize, text=bkGreen!70!black, anchor=south] at (29,7.5) {\qty{25}{m/s}};
  \draw[bk arrow, bkPurple] (92,6) -- (102,6);
  \node[font=\sffamily\footnotesize, text=bkPurple!80!black, anchor=south] at (97,7.5) {\qty{15}{m/s}};
  \draw[{Stealth[length=4pt]}-{Stealth[length=4pt]}, bkInk, line width=0.7pt] (19,-5.5) -- (68,-5.5);
  \node[font=\sffamily\footnotesize, text=bkInk, anchor=north] at (43.5,-6.5) {\qty{40}{m}};
\end{tikzpicture}
\end{bkfigure}

\bkproblem あめを\R{1人}{ひとり}3\R{個}{こ}ずつ\R{配}{くば}ると8\R{個}{こ}あまり、5\R{個}{こ}ずつ\R{配}{くば}ると6\R{個}{こ}たりません。\R{人数}{にん|ずう}を $x$ \R{人}{にん}として\R{方程式}{ほう|てい|しき}をつくり、\R{人数}{にん|ずう}とあめの\R{数}{かず}を\R{求}{もと}めましょう。
\bkpans{$3x+8=5x-6$、$-2x=-14$、$x=7$。7\nobreak\R{人}{にん}、あめは $3\times7+8=29$ で29\nobreak\R{個}{こ}}

\bkproblem \R{兄}{あに}は\R{秒速}{びょう|そく}\qty{3}{m}、\qty{10}{m}\R{先}{さき}の\R{弟}{おとうと}も\R{秒速}{びょう|そく}\qty{3}{m}で\R{同}{おな}じときに\R{走}{はし}り\R{出}{だ}しました。$3t=10+3t$ を\R{解}{と}こうとすると、どうなりますか。それは\R{場面}{ば|めん}では\R{何}{なに}を\R{表}{あらわ}していますか。
\bkpans{\R{移項}{い|こう}すると $0=10$ となり、\R{解}{かい}がない。いつまでも\qty{10}{m}はなれたままで、\R{追}{お}いつかない}

\begin{furikaeri}
\begin{checklist}
\item \R{方程式}{ほう|てい|しき}をつくり、\R{表}{ひょう}でためして\R{解}{かい}をさがせる。\hfill{\footnotesize\sffamily\color{bkTheme}→\pageref{bk:2-16:houteishiki}ページ}
\item 1\R{次方程式}{じ|ほう|てい|しき}を1\R{行}{ぎょう}ずつ\R{解}{と}ける。\hfill{\footnotesize\sffamily\color{bkTheme}→\pageref{bk:2-16:tejun}ページ}
\item \R{表}{ひょう}・グラフ・\R{方程式}{ほう|てい|しき}で\R{解}{と}き、\R{確}{たし}かめられる。\hfill{\footnotesize\sffamily\color{bkTheme}→\pageref{bk:2-16:graph}ページ}
\end{checklist}
\end{furikaeri}

\bkfukidashi{sora}{\R{表}{ひょう}でもグラフでも\R{方程式}{ほう|てい|しき}でも、\R{答}{こた}えは5\R{秒}{びょう}。\R{方程式}{ほう|てい|しき}なら、4.5\R{秒}{びょう}のような\R{答}{こた}えもぴったり\R{出}{で}るね。}
